\section{Reconstruction, native constraints and parameter coverage}\label{sec:reconstruction}
This section supplies the intermediate steps behind Main Sections~3.1--3.3.
The battery energy storage system (BESS) query retains the complete discharge
positive net power trajectory and any unchanged variables. Stored energy is
denoted by $E$, power by $c,d,p$, and the known energy offset by $b$.
Equal query sets are abbreviated EQ; STRICT denotes strict inclusion of the reference query set in the relaxed query set.
State of charge (SOC) is a dimensionless fraction or a percentage; it is not
interchanged with stored energy without a documented energy scale.

\subsection{The structural candidate and the role of uniqueness}
At a fixed parameter state, the reconstruction is defined in the structural
set before its operating feasibility is tested. Nonnegative directional flows
and exclusivity fix the charge and discharge powers from net power. Known
initial energy and the deterministic recursion then fix every subsequent
energy. Unchanged query variables supply the remaining coordinates.
The operating inequalities determine whether this candidate is feasible.

The reverse implication in Main Proposition~1 uses the fact that each
structural fiber is a singleton. If an admitted query has a feasible reference
realization, that realization must be the same structural candidate. A
nonunique fiber would only establish existence of a feasible realization.
The following counterexample isolates this role of uniqueness:
\begin{equation}
\mathcal P=\{a, b\},\quad R=X=\{a\},\quad
\pi(a)=\pi(b)=0,\quad L(0)=b .\label{eq:nonunique}
\end{equation}
The two projected sets coincide, but this particular selector is infeasible.
Sufficiency still holds for any feasible query preserving selector. The
battery specialization avoids this ambiguity through its unique lift.


Table~\ref{tab:assumption_roles} distinguishes the premises of the set argument,
the battery reconstruction and the finite computational criterion. Coefficients
may vary across intervals while remaining fixed with respect to the assessed
parameter state. Each interval has one charge and discharge pair in the
declared model. Allowing unrestricted switching within an interval would
change the reference trajectory space and the reconstruction question.

\begin{table}[htbp]
\caption{Distinct roles of the mathematical premises.}\label{tab:assumption_roles}
\begin{tabularx}{\linewidth}{@{}p{.22\linewidth}p{.41\linewidth}X@{}}
\toprule Argument & Required structure & Consequence\\\midrule
Projected set equivalence & Containment; the same query on both models; a total, unique structural lift for each admitted query & Identifies any feasible reference realization with that lift. Affinity and compactness are not premises of this set argument.\\[3pt]
Battery reconstruction & Known initial energy and coefficients; finite horizon; nonnegative directional powers; positive efficiencies & Exclusivity fixes both powers and the deterministic recursion fixes energy. Unknown initial states or nonunique aggregate realizations require different quantifiers.\\[3pt]
Energy ordering & Positive durations, nonnegative retention and efficiency product at most one & The energy shift is nonnegative. These sign conditions are distinct from the algebraic reconstruction identity and projected set logic.\\[3pt]
Attained native residual maximum & Nonempty compact relaxed set; continuous reconstruction; complete finite affine operating rows & Every required inequality and both directions of each equality are preserved exactly when all row maxima are nonpositive.\\[3pt]
Linear parameter characterization & Fixed horizon and coefficients; affine initial energy, offsets and right hand sides & Each row and sign region gives a linear parameter system. Variable coefficients, uncertain dynamics and invariance do not follow from this conclusion.\\
\bottomrule\end{tabularx}
\end{table}

\subsection{Energy propagation and dimensional consistency}
Let $w_t=\min(c_t,d_t)$. Substituting the absolute value forms of the
exclusive powers into the energy input gives the intermediate identities
\begin{equation}\begin{aligned}
\eta_{c, t}c_t^*-d_t^*/\eta_{d, t}
 &=-\alpha_t p_t-\tfrac 12\kappa_t|p_t|,\label{eq:input_star}\end{aligned}\end{equation}
\begin{equation}\begin{aligned}
\eta_{c, t}c_t-d_t/\eta_{d, t}
 &=-\alpha_t p_t-\tfrac 12\kappa_t|p_t|-\kappa_tw_t .\label{eq:input_relaxed}
\end{aligned}\end{equation}
Here $\kappa_t=1/\eta_{d,t}-\eta_{c,t}$ and
$\alpha_t=(\eta_{c,t}+1/\eta_{d,t})/2$, as in Main Section~3.2.
Subtracting the two recursions eliminates the common known offset:
\begin{equation}
\delta E_t=\rho_t\delta E_{t-1}+h_t\kappa_tw_t,
\qquad \delta E_0=0 .\label{eq:shift_recurrence}
\end{equation}
Forward substitution assigns the input at interval $t$ to each later
energy through the product of subsequent retention factors. This produces
the propagation matrix $G$ in the Main, with $G_{tt}=h_t$ and an empty
retention product equal to one. Stacking the intervals yields
\begin{equation}
\delta E=G(\kappa\odot w).\label{eq:shift_stack}
\end{equation}
The offset is already an energy quantity and is not multiplied by duration
a second time. The matrix has time units, the simultaneous component has
power units, and their propagated product has energy units. Positive
durations, nonnegative retention and an efficiency product at most one
make the energy shift componentwise nonnegative. These sign conditions
explain energy ordering; they do not establish feasibility of every native row.

\subsection{Completeness and preservation of affine constraints}
The constraint index covers every native operating inequality whose
satisfaction must survive reconstruction. Deterministic dynamics already
enforced by the structural lift are not duplicated. A retained operating
equality contributes both inequality directions. Couplings involving the
unchanged variables remain in the same row system.

For a native row, the power terms at the reconstruction expand as
\begin{equation}
a_c^\top c^*+a_d^\top d^*
=\tfrac 12(a_d-a_c)^\top p+\tfrac 12(a_c+a_d)^\top|p| .
\label{eq:power_row}
\end{equation}
Combining this expansion with energy propagation gives the coefficients
$s_j,\ell_j,v_j$ defined in Main Section~3.3. The row change is
\begin{align}
\Delta_j
&=a_{E, j}^\top G(\kappa\odot w)
 -(a_{c, j}+a_{d, j})^\top w\\
&=-2v_j^\top w .\label{eq:row_change}
\end{align}
The unchanged variables cancel from this difference. If $v_j$ is
componentwise nonnegative, the row cannot increase. A negative component
does not itself prove a violation: feasible simultaneous flows and the
original row slack still matter. Such rows require evaluation over the
complete relaxed feasible set.

The maximum residual in the Main is attained because the relaxed set is
nonempty and compact and the reconstructed row is continuous. Its sign,
rather than agreement on a selected dispatch, determines preservation of
that row for all queries. Both directions of equalities must pass. If a row
has units $U_j$, its residual and row change retain $U_j$. Residuals with
different physical units or row scalings are not combined into a distance,
probability or common physical margin.

\subsection{Parameter regions and the scope of a reuse certificate}
The pointwise criterion uses fixed known coefficients during an assessment.
For a parameter region, keep horizon, durations, efficiencies, retention and
native row coefficients fixed, while initial energy, known offsets and
right hand sides may vary affinely. The zero power energy trajectory is
then affine in the parameters. Within each net power sign region, relaxed
feasibility and a positive reconstructed residual form linear systems in
the trajectory and parameter variables. Projecting these systems and taking
their finite union identifies the parameters admitting a violating query.
The explicit sign-region reconstruction is
\begin{equation}
E^*_{\sigma}(\theta,p)=f(\theta)-G\left(\alpha\odot p+
\frac{\kappa\odot\sigma\odot p}{2}\right),\qquad
\sigma\in\{-1,1\}^T,\quad \sigma\odot p\geq0.
\label{eq:parameter_sign_lift}
\end{equation}
For each row and sign vector, define the joint parameter and trajectory set
\begin{equation}
\begin{split}
\mathcal B_{j,\sigma}=\bigl\{(\theta,y):\;&y\in R(\theta),\quad
p=d-c,\quad\sigma\odot p\geq0,\\
&a_{E,j}^{\top}f(\theta)+\ell_j^{\top}p+
v_j^{\top}(\sigma\odot p)+a_{z,j}^{\top}z>r_j(\theta)\bigr\}.
\end{split}\label{eq:parameter_bad_set}
\end{equation}
The trajectory and parameter dependence is affine because $G$, $\alpha$,
$\kappa$, $\ell_j$ and $v_j$ are fixed with respect to $\theta$.
Let $\operatorname{proj}_{\theta}$ eliminate the trajectory coordinates.
The nonempty and valid parameter sets are therefore
\begin{equation}
D=\operatorname{proj}_{\theta}\{(\theta,y):y\in R(\theta)\},\qquad
V=D\setminus\bigcup_{j,\sigma}\operatorname{proj}_{\theta}\mathcal B_{j,\sigma}.
\label{eq:parameter_region_projection}
\end{equation}
Finite elimination of linear equalities and strict or nonstrict inequalities
gives a finite Boolean linear description. Zero powers can lie in both
adjacent sign regions because their reconstructed values agree. Strict
inequalities preserve the difference between a violating query and an
admissible boundary; closing the violation sets would change that meaning.
This construction need not yield a small or convex region.

Distinguish the true validity region $V$ from a proved
covered set $\Omega\subseteq V$. For a feasible parameter state,
\begin{equation}
\theta\in\Omega\ \Longrightarrow\ \mathcal Q_R(\theta)=\mathcal Q_X(\theta),
\qquad
\theta\in D\setminus V\ \Longrightarrow\
\mathcal Q_X(\theta)\subsetneq\mathcal Q_R(\theta).
\label{eq:coverage}
\end{equation}
Leaving $\Omega$ alone establishes neither relation.





\begin{table}[htbp]

\caption{Scope of reuse after changes to a certified storage model.}
\label{tab:si_certificate_reuse}

\begin{tabularx}{\linewidth}{@{}p{.25\linewidth}XX@{}}
\toprule
Change & Available guarantee & Condition for reuse\\
\midrule
Initial energy, known offsets or operating limits
& Equality throughout the proved parameter coverage
& The complete updated state remains inside that coverage with its fixed coefficients.\\
Efficiencies, retention or native row coefficients
& The pointwise criterion remains applicable if its assumptions hold
& Assess the revised model; the earlier parameter region does not automatically extend.\\
Horizon or interval definition
& The criterion can assess the revised finite trajectory space
& Establish coverage for the new dynamics and complete query.\\
Common restriction on certified query outputs
& Equality is preserved by the same restriction in both query spaces
& Restrict only outputs already retained by the certificate.\\
Native operating constraints
& The modified model pair can be assessed
& Verify containment, unique reconstruction and all revised native rows.\\
Reference model or query outputs
& A common mapping of certified outputs preserves their set equality for the same reference
& A changed reference or outputs not determined by the certified query require separate justification.\\
\bottomrule
\end{tabularx}
\end{table}

An identical restriction on certified queries differs from imposing an
additional condition on internal charge, discharge or energy variables.
The latter condition requires reconstruction feasibility to be checked
again. Equality of restricted query sets can include an empty restricted
set; operational use additionally requires an admissible request.
Coverage does not establish that future operation remains inside the
covered region.


\section{Two interval validity boundaries}\label{sec:boundaries}
The BOX, END and START models in Main Section~4.1 use the same energy
scale and dynamics. This analytical study is separate from the illustrative
device used for the state transitions in Section~\ref{sec:transitions}.
There are two intervals of $0.1$~h, charging and discharging efficiencies
of $0.9$, unit retention and zero energy offset. Energy bounds are
$[0,1]$~MWh, the shared power bound is $c_t+d_t\leq1$~MW, and terminal
energy is free within the energy bounds. The initial energy is fixed in
each assessment.

BOX retains only these power and energy bounds. The coupled models add
a charging ceiling evaluated at different interval endpoints:
\begin{equation}
\begin{array}{ll}
\mathrm{END}:&c_t+\beta E_t\leq A,\\
\mathrm{START}:&c_t+\beta E_{t-1}\leq A,
\end{array}
\qquad \beta=1~\mathrm{h}^{-1}.\label{eq:boundary_models}
\end{equation}
The ceiling $A$ has power units. The coupling acts on charging power,
not on discharge positive net power. The numerical parameter domain is
$0\leq E_0\leq1$~MWh and $0.2\leq A\leq2$~MW.

\subsection{Exact feasible and valid domains}
Use dimensionless coordinates $e=E_0/(1~\mathrm{MWh})$ and
$a=A/(1~\mathrm{MW})$. Define $I(u)=[0,u]$ for $u\geq0$ and
$I(u)=\varnothing$ otherwise. Table~\ref{tab:boundary_domains} gives
the exact domains in these coordinates.

\begin{table}[htbp]
\caption{Feasible and valid initial energy domains of the two interval models.}
\label{tab:boundary_domains}
\begin{tabularx}{\linewidth}{@{}l p{.25\linewidth} X@{}}
\toprule Model & Nonempty domain $D$ & Validity domain $V$\\\midrule
BOX & $[0,1]$ & $[0,41/50]$\\[4pt]
END & $[0,\min(1,a+1/9)]$ &
$\bigl(\{a+1/9\}\cap[0,1]\bigr)\cup I(\min(41/50,a-59/50))$\\[4pt]
START & $[0,\min(1,a)]$ &
$\bigl(\{a\}\cap[0,1]\bigr)\cup I(\min(41/50,a-109/100))$\\
\bottomrule\end{tabularx}
\end{table}

The singleton terms retain the one dimensional valid branches at the lower
feasibility boundaries. A state in $D\setminus V$ has strict query set
inclusion. A state outside $D$ has an empty relaxed model and is not a
validity certificate at an admissible operating state.

For the two dimensional valid portions, the charging ceilings differ by
$59/50-109/100=0.09$ in the fixed numerical power coordinate.
At $E_0=0.30$~MWh and $A=1.40$~MW, BOX and START are valid, while
END is feasible with strict inclusion. These statements include the exact
boundaries rather than an interpolation through plotted cells.

\subsection{Recorded quantified checks and discrete verification}
The saved decision systems compare each predicted domain with the native
model in both directions. For each model they test feasibility outside the
proposed nonempty domain, emptiness inside that domain, a violating query
inside the proposed valid domain, and a valid state outside that domain.
All four violation systems were recorded as unsatisfiable for each of the
three models. The corresponding quantified formulas and results are
retained with the computational material. Quantifier elimination and its
exact formula records provide a separate source of the boundary expressions.

Table~\ref{tab:boundary_grid} summarizes the accompanying grid of 330 fixed checks, 110 per model. The
selected coordinates are $e=0,0.1,\ldots,1$ and
$a=0.2,0.4,\ldots,2$. They verify the stated classifications at those
coordinates. The continuous characterization follows from the recorded
quantified checks, not from the number or visual density of grid cells.


\begin{table}[htbp]
\caption{Exact classifications at the 330 prescribed model and parameter states.}
\label{tab:boundary_grid}
\begin{tabular}{@{}lrrrr@{}}\toprule
Model & EQ & STRICT & Empty & Total\\\midrule
BOX &90&20&0&110\\
END &25&69&16&110\\
START &34&56&20&110\\
\bottomrule\end{tabular}
\end{table}


\subsection{Logical form of the continuous boundary checks}
The twelve saved decisions concern the declared parameter rectangle, not
just the grid. Let $P(e,a)$ denote that rectangle, let $F(e,a,y)$ denote
native relaxed feasibility, and let $B(e,a,y)$ denote relaxed feasibility
together with at least one strict reconstructed violation. Write
$\widehat D$ and $\widehat V$ for the expressions in
Table~\ref{tab:boundary_domains}. For each model the four counterexample
systems are
\begin{equation}\begin{aligned}
&\exists e,a,y:\ P(e,a)\land\neg\widehat D(e,a)\land F(e,a,y),
\label{eq:domain_check_one}\end{aligned}\end{equation}
\begin{equation}\begin{aligned}
&\exists e,a:\ P(e,a)\land\widehat D(e,a)\land
                 \forall y\,\neg F(e,a,y),\label{eq:domain_check_two}\end{aligned}\end{equation}
\begin{equation}\begin{aligned}
&\exists e,a,y:\ P(e,a)\land\widehat V(e,a)\land B(e,a,y),
\label{eq:domain_check_three}\end{aligned}\end{equation}
\begin{equation}\begin{aligned}
&\exists e,a:\ P(e,a)\land\widehat D(e,a)\land\neg\widehat V(e,a)
                 \land\forall y\,\neg B(e,a,y).\label{eq:domain_check_four}
\end{aligned}\end{equation}
Their recorded unsatisfiability establishes both inclusions of each proposed
domain. The universal variables range over the complete relaxed flow vector;
feasibility is included in $F$ and $B$. Strict violations retain admissible
boundary equalities. Four charge or discharge sign assignments make each
reconstruction affine, including zero power in either adjacent region.
The companion \texttt{data/Boundary\_Check\_Specification.json} states the exact
rational recursions and each native row. The original builder and closed
region script are supplied with their original module names under
\texttt{code/original\_layout/RESEARCH/round3/}. Earlier incomplete quantifier
elimination records are not equality certificates. The saved decisions are
not newly solved by the saved record verifier.


\section{State transition design and saved reconstruction witnesses}\label{sec:transitions}
\subsection{Illustrative device and native operating constraints}
The device settings follow Rosewater's illustrative energy reservoir model,
Section~2.3.1, Eqs.~(2.3.4) and (2.3.6)--(2.3.9), and Table~2.5
\cite{rosewater2020battery}. The source represents an illustrative device,
rather than an identified installation. Its energy coordinate places
conversion loss on charging. The source charging efficiency is $0.65$;
$0.90$ is a sensitivity setting and $1$ is a lossless control. Discharging
efficiency and retention are one. The common parameters and their roles
are consolidated in Table~\ref{tab:internal_parameters}.

\begin{table}[htbp]
\caption{Illustrative device parameters and selected study settings.}\label{tab:internal_parameters}
\begin{tabularx}{\linewidth}{@{}p{.33\linewidth} p{.26\linewidth} X@{}}
\toprule Quantity & Value & Origin or use\\\midrule
Nominal energy capacity & $0.6$ MWh & Source SOC to energy scale\\
SOC limits / source initial SOC & $20$ to $95\%$ / $60\%$ & Source fractions of nominal capacity\\
Energy bounds & $0.12$ to $0.57$ MWh & Source capacity and SOC range\\
Shared power rating & $0.5$ MW & Study split power bound\\
Charging taper coefficient & $50/3~\mathrm{h}^{-1}$ & Source signed power constraint\\
Discharging taper coefficient & $25/3~\mathrm{h}^{-1}$ & Source signed power constraint\\
Auxiliary power & $0.007$ MW & Constant energy drain\\
Charging efficiency & $0.65,0.90,1$ & Source, sensitivity and control\\
Source initial energy & $0.36$ MWh & Source initial SOC\\
Charging taper knee & $0.54$ MWh & Boundary initialization reference\\
Initialization margin & $0.003$ MWh & Selected study margin\\
Interval duration & $1,5,15$ min & Selected temporal discretizations\\
\bottomrule\end{tabularx}
\end{table}

The source's printed charge/discharge labels and footnote are inconsistent
with the adjacent signed equation. The mapping follows the signed equation
and its explicit slope and intercept expressions. The shared split power
bound and enforcement at both interval endpoints are declared study choices.
The source economic objective and terminal cycle constraint are not retained.

Both relaxed formulations satisfy
\begin{equation}\begin{aligned}
E_t&=E_{t-1}+h(\eta_c c_t-d_t-P_{\mathrm{aux}}),\label{eq:internal_dynamics}\end{aligned}\end{equation}
\begin{equation}\begin{aligned}
c_t, d_t&\geq 0,\qquad c_t+d_t\leq P,\qquad
 E_{\min}\leq E_t\leq E_{\max}.\label{eq:internal_box}
\end{aligned}\end{equation}
For both $s=t-1$ and $s=t$, the net power relaxation also imposes
\begin{equation}
c_t-d_t+\beta_cE_s\leq\beta_cE_{\max},\qquad
d_t-c_t-\beta_dE_s\leq-\beta_dE_{\min}.\label{eq:internal_net}
\end{equation}
The gross power tightening replaces the flow terms by $c_t$ and $d_t$:
\begin{equation}
c_t+\beta_cE_s\leq\beta_cE_{\max},\qquad
d_t-\beta_dE_s\leq-\beta_dE_{\min}.\label{eq:internal_gross}
\end{equation}
The exclusive reference adds $c_td_t=0$. Under the stated energy bounds,
the two formulations represent the same exclusive device, but their
relaxed feasible sets differ. In each assessment the query is the entire
net power trajectory $p=d-c$.

\subsection{Initialization and feasibility of the state update}
The source initialization fixes $E_0=0.36$~MWh. The second initialization
places the state relative to the charging boundary:
\begin{equation}
E_0=0.54~\mathrm{MWh}
-Th\bigl(\eta_c\,0.5~\mathrm{MW}-0.007~\mathrm{MW}\bigr)
-0.003~\mathrm{MWh}.\label{eq:initialization}
\end{equation}
Here $h$ is expressed in hours. One exclusive full power charging action
then gives the updated energy
\begin{equation}
E_1=E_0+h(\eta_cP-P_{\mathrm{aux}}).\label{eq:update}
\end{equation}
The two assessments use $E_0$ and $E_1$ as their respective fixed initial
energies while retaining horizon, coefficients, constraints and query.
The update action is separate from the assessment horizon.

Eligibility requires both initial energies and the actual update action to
satisfy the rating, energy bounds and both endpoint native power limits.
Invalid inputs are not clipped or replaced. Both slots of such a pair remain
in the design as inapplicable. They are neither empty relaxed models nor
unresolved solver assessments.

\subsection{Factor design and accounting units}
At each selected horizon through 24 intervals, the design crosses three
charging efficiencies, three durations, two row formulations and two
initializations. This gives 36 pairs and 72 state slots per horizon.
The evaluated horizons are $4,5,6,7,8,10,12,14,16,20,24$.
At $32,48,96,128,192,256$ intervals, the panel retains all efficiencies
and both formulations, but only the one and fifteen minute durations and
the source initialization. This gives 12 pairs and 24 state slots per horizon.
The conditional entry rules are stated in Section~\ref{sec:computation}.

The scientific design therefore contains 468 pairs and 936 state slots.
The separate two interval regression has 36 pairs and 72 slots and is used
only to compare implementations. It is excluded from the scientific totals.
Table~\ref{tab:horizons} preserves the horizon specific denominators and
lossless controls. These are structured model assessments, not a sample
of independent devices or an estimate of operating failure frequency.
\begin{table}[htbp]
\begingroup
\caption{State assessments and lossless controls at the evaluated horizons.}\label{tab:horizons}
\begin{tabular}{@{}rrrrrrrrr@{}}
\toprule
 &\multicolumn{5}{c}{All state assessments}&\multicolumn{3}{c}{Lossless subset}\\

$T$ & States & EQ & STRICT & Unresolved & Inapplicable & EQ & Unresolved & Inapplicable\\
\midrule
4 & 72 & 44 & 20 & 0 & 8 & 20 & 0 & 4\\
5 & 72 & 42 & 22 & 0 & 8 & 20 & 0 & 4\\
6 & 72 & 38 & 22 & 0 & 12 & 20 & 0 & 4\\
7 & 72 & 36 & 24 & 0 & 12 & 20 & 0 & 4\\
8 & 72 & 36 & 24 & 0 & 12 & 20 & 0 & 4\\
10 & 72 & 36 & 24 & 0 & 12 & 20 & 0 & 4\\
12 & 72 & 30 & 22 & 0 & 20 & 16 & 0 & 8\\
14 & 72 & 15 & 22 & 15 & 20 & 8 & 8 & 8\\
16 & 72 & 0 & 20 & 28 & 24 & 0 & 16 & 8\\
20 & 72 & 0 & 20 & 28 & 24 & 0 & 16 & 8\\
24 & 72 & 0 & 22 & 26 & 24 & 0 & 16 & 8\\
32 & 24 & 0 & 12 & 12 & 0 & 0 & 8 & 0\\
48 & 24 & 0 & 16 & 8 & 0 & 0 & 8 & 0\\
96 & 24 & 0 & 16 & 8 & 0 & 0 & 8 & 0\\
128 & 24 & 0 & 16 & 8 & 0 & 0 & 8 & 0\\
192 & 24 & 0 & 14 & 10 & 0 & 0 & 8 & 0\\
256 & 24 & 0 & 15 & 9 & 0 & 0 & 8 & 0\\
\bottomrule\end{tabular}
\par\smallskip\raggedright Each state is a fixed parameter assessment. Lossless controls are subsets of the same design. The separate two interval implementation regression is excluded. No lossless state was certified STRICT.
\endgroup
\end{table}


\subsection{Conditional composition of the paired outcomes}
The paired classifications are grouped by charging efficiency and initialization
in Table~\ref{tab:residual_paired_conditions}. The source initialization fixes
initial stored energy at $0.36$~MWh; the boundary relative initialization uses
Eq.~\eqref{eq:initialization}. Each entry counts
a predefined pair of before and after validity assessments with the same horizon,
coefficients, operating constraints and complete net power query.

\begin{table}[htbp]
\caption{Paired outcome conditions.}\label{tab:residual_paired_conditions}

\begin{tabular}{@{}c l r r r r r r@{}}
\toprule
$\eta_c$ & Initialization & Pairs & EQ$\to$EQ & EQ$\to$STRICT & STRICT$\to$STRICT & Unresolved pair & Inapplicable pair\\
\midrule
0.65 & Source & 90 & 19 & 2 & 59 & 10 & 0 \\
0.65 & Boundary relative & 66 & 0 & 33 & 0 & 9 & 24 \\
0.90 & Source & 90 & 15 & 2 & 64 & 9 & 0 \\
0.90 & Boundary relative & 66 & 0 & 27 & 0 & 7 & 32 \\
1.00 & Source & 90 & 44 & 0 & 0 & 46 & 0 \\
1.00 & Boundary relative & 66 & 28 & 0 & 0 & 6 & 32 \\
\midrule
 & Total & 468 & 106 & 64 & 123 & 87 & 88\\
\bottomrule
\end{tabular}
\par\smallskip
\begin{minipage}{\linewidth}
EQ denotes complete query set equality, and STRICT denotes strict inclusion.
An unresolved pair has at least one assessment without an accepted decision;
the other assessment may be certified. Inapplicable pairs fail the prescribed
initialization or update eligibility checks. The restricted panel above
24 intervals includes only the source initialization, so the row totals do
not represent matched sampling across the two initializations.
\end{minipage}
\end{table}

Of the 64 certified equality to strict inclusion transitions, 60 used the
boundary relative initialization and four used the source initialization.
All occurred at evaluated horizons from 4 to 14 intervals with charging
efficiency below one. At every transition, the retained energy admissible
reconstruction witnesses violated the charging constraint at the end of the
last interval. These observations describe the saved witness set; they do
not establish uniqueness of the violated constraint across all possible
queries or absence of transitions in unresolved settings.

Matching the two relaxed row formulations at the same parameter state gave
eight differing classification labels, each involving an unresolved
assessment. There was no disagreement between two certified labels.
Agreement of classification labels does not assert equality of the relaxed
query sets of the two formulations.


\subsection{Representative native power violation}
The four interval net power example uses one minute intervals and
$\eta_c=0.65$. The eligible charging update changes initial energy from
$0.5158$ to $0.5211$~MWh. Table~\ref{tab:transition_vector} records the
saved relaxed trajectory and its exclusive reconstruction at the updated
state. Discharging after reconstruction is zero, so the common query is
$p_t=-c_t^*$. The charging row residual is
\begin{equation}
r_t=c_t^*-d_t^*+\frac{50}{3}E_t^*-\frac{19}{2} .\label{eq:witness_residual}
\end{equation}
Power is in MW and energy in MWh. The last residual is exactly
$161/19800$~MW. The complete relaxed trajectory is feasible; reconstructed
energy remains between $0.12$ and $0.57$~MWh. All other native inequalities
hold for this witness. A saved excess is distinct from a maximum over all
queries, even when another verification route reports the same value.

\begin{table}[htbp]
\caption{Saved four interval trajectory and exclusive reconstruction.}\label{tab:transition_vector}
\begin{tabular}{@{}rrrrrrr@{}}\toprule
$t$ & $c_t$ & $d_t$ & $c_t^*$ & $E_t$ & $E_t^*$ & $1000r_t$\\
 & MW & MW & MW & MWh & MWh & kW\\\midrule
1&0.416364&0.083636&0.332727&0.524100&0.524588&$-424.141414$\\
2&0.500000&0&0.500000&0.529400&0.529888&$-168.535354$\\
3&0.500000&0&0.500000&0.534700&0.535188&$-80.202020$\\
4&0.500000&0&0.500000&0.540000&0.540488&$8.131313$\\
\bottomrule\end{tabular}
\end{table}

The table is rounded for display. Exact fractional vectors and all native
row residuals are retained in the accompanying computational records.
Each of the original 64 certified transitions is linked to its two assessed states and at least
one saved energy feasible reconstruction with a native power violation.
This existence statement does not describe every admitted relaxed query.

\subsection{Earlier two interval controls}
Table~\ref{tab:t2_controls} retains nine settings from the earlier two interval
benchmark represented in the implementation regression. The boundary relative
initialization changes with both efficiency and duration; the updated energy
is therefore stated for each setting. The power values are positive parts of
saved numerical LP maxima subject to feasible reconstructed energy. Stored
exact logical queries separately test the existence of an energy feasible
native power violation. These are prescribed comparisons at different initial
states, rather than an isolated duration effect at a common state. A zero
numerical maximum alone does not establish complete query equality.

\begin{table}[H]
\caption{Saved two interval controls under boundary relative initialization.}
\label{tab:t2_controls}
\begin{tabular}{@{}rrrrl@{}}\toprule
$\eta_c$ & Duration (min) & Updated $E_0$ (MWh) & Excess (kW) & Exact existence query\\\midrule
0.65 & 1 & 0.531700 & 8.131 & Satisfiable\\
0.65 & 5 & 0.510500 & 83.081 & Satisfiable\\
0.65 & 15 & 0.457500 & 270.455 & Satisfiable\\
0.90 & 1 & 0.529617 & 3.845 & Satisfiable\\
0.90 & 5 & 0.500083 & 29.751 & Satisfiable\\
0.90 & 15 & 0.426250 & 94.518 & Satisfiable\\
1.00 & 1 & 0.528783 & 0.000 & Unsatisfiable\\
1.00 & 5 & 0.495917 & 0.000 & Unsatisfiable\\
1.00 & 15 & 0.413750 & 0.000 & Unsatisfiable\\
\bottomrule\end{tabular}
\end{table}

The lookup and original records are linked to the generalized regression by
their actual input parameters. They remain part of the separate two interval
implementation evidence and are not added to the 936 scientific state slots
or the 64 certified transitions. The LP maxima in this table are distinct
from the selected four interval witness residual in Table~\ref{tab:transition_vector}.


\section{Computational certification and timing}\label{sec:computation}


Classifications and computation times in this section describe the original
verification protocol. Section~\ref{sec:sat_certificate_followup} reports the
subsequent certificates for the same inputs, including the 104 originally
unresolved lossless cases. The retained figures and tables preserve the
original protocol outcomes.

\subsection{Reduced linear programming evaluation}
Linear programming (LP) evaluates the reconstructed native row after
substitution of the exclusive lift. The finite formulation below applies to
the continuous affine relaxed models used in the computational study.
Rows shown to be automatically preserved in Section~\ref{sec:reconstruction}
remain native constraints even when their optimization is omitted.

For a row requiring evaluation, define $J_j^+=\{t:v_{j,t}>0\}$ and
$J_j^-=\{t:v_{j,t}<0\}$. For each assignment
$\sigma\in\{-1,1\}^{J_j^+}$, the reduced subproblem is
\begin{equation}
\begin{aligned}
\max_{y, u}\quad & a_{E, j}^\top f+\ell_j^\top p+a_{z, j}^\top z-r_j
 +\sum_{t\in J_j^+}v_{j, t}\sigma_tp_t
 +\sum_{t\in J_j^-}v_{j, t}u_t\\
\text{subject to}\quad &y\in R(\theta),\quad p=d-c,\\
&u_t\geq p_t,\quad u_t\geq-p_t,\qquad t\in J_j^-.
\end{aligned}\label{eq:reduced_lp}
\end{equation}
Negative coefficients make the epigraph tight at an optimum. Positive
absolute value terms are represented by the finite maximum over signs.
Taking the largest subproblem optimum therefore gives the mathematical row
maximum. This reduced formulation does not impose sign restrictions on
the net power itself. A positive absolute value coefficient cannot instead
multiply an unconstrained epigraph variable in this maximization.

The planned computational workload is one relaxed feasibility LP plus
$\sum_j2^{|J_j^+|}$ for the rows selected for optimization. If a completed
positive maximum dominates the zero upper bounds for automatically
preserved rows, it is also the maximum over the target rows. A nonpositive
value computed only over selected rows is not their signed global maximum.
A best value at interruption is a lower bound on an unfinished maximum.

\subsubsection*{Simpler rows and fixed sign queries}
The enumeration bound is sufficient for the stated implementation, not a
lower bound on certification cost. A reconstructed charging ceiling has the
elementary equivalent form
\begin{equation}
\max\{-p_t,0\}+\beta E_t^*\leq A
\quad\Longleftrightarrow\quad
\begin{cases}
-p_t+\beta E_t^*\leq A,\\
\beta E_t^*\leq A.
\end{cases}\label{eq:charging_normalization}
\end{equation}
For $\beta\geq0$, the second inequality is automatic if reconstructed energy
is already known to obey an upper bound $U$ throughout the assessed relaxed
set and $A\geq\beta U$. Under the propagation and efficiency sign conditions
of Section~\ref{sec:reconstruction}, $E_t^*$ is concave in net power, so the
remaining row maximum has an ordinary LP representation using negative
absolute value coefficients. The reconstructed energy premise must be
established separately; it cannot be inferred from relaxed energy feasibility.
For $\beta<0$, cancellation instead decreases the original charging row
directly. A row $d_t-\beta E_t\leq A$ is similarly preserved for $\beta\geq0$.

If all net power signs are fixed by the assessed operating envelope, the
entire reconstruction is affine and ordinary polyhedral row implication
suffices. These elementary special cases are retained as baselines for
interpreting the general signed row formulation. They add no solver runs
and do not change the recorded implementation counts or timings.

\subsection{Full sign reference and exact logical decision}
The full sign reference enumerates $\sigma\in\{-1,1\}^T$, imposes
$\sigma_tp_t\geq0$, and substitutes $|p_t|=\sigma_tp_t$. Infeasible sign
regions are omitted. Zero powers can belong to adjacent regions without
changing the maximum. In the horizon study this reference is used only for
$T\leq8$. It also solves the existing constrained power violation problem
with reconstructed energy bounds imposed. Its workload therefore differs
from that of the reduced formulation.

Satisfiability modulo theories (SMT) supplies an exact route over linear
real arithmetic. At fixed rational coefficients, conditional linear
expressions encode the max formulas for exclusive charge and discharge.
The native relaxed constraints and reconstructed energy recursion remain
explicit. The decision first tests relaxed nonemptiness, then the existence
of a violating relaxed query:
\begin{equation}
\exists y\in R(\theta):\quad
\bigvee_{j\in\mathcal J}
\left[a_{E, j}^\top E^*+a_{c, j}^\top c^*
+a_{d, j}^\top d^*+a_{z, j}^\top z>r_j\right].\label{eq:smt_violation}
\end{equation}
Equalities contribute both directions. Unsatisfiability of this violation
query, together with nonemptiness, certifies complete query equality.
A satisfiable result supplies a candidate whose native feasibility and
violation are rechecked exactly. A separate diagnostic query asks whether
reconstructed energy can remain feasible while a native power row fails.
This is an existence test, not a maximum power violation computation.

\subsection{Certificates and saved witnesses}
Table~\ref{tab:semantics} distinguishes certified query relations from unavailable assessments.

\begin{table}[htbp]
\caption{Interpretation of state assessments and unavailable decisions.}\label{tab:semantics}
\begin{tabularx}{\linewidth}{@{}p{.20\linewidth} X@{}}\toprule
Outcome & Required evidence\\\midrule
Equality & A nonempty relaxed model and exact absence of a violating query\\
Strict inclusion & At least one exactly verified relaxed feasible query with an infeasible exclusive reconstruction\\
Empty model & Exact relaxed infeasibility, placing the parameter outside the declared nonempty domain\\
Unresolved & Neither accepted exact relation is established within the permitted computation\\
Inapplicable & The designed initialization or update fails the specified eligibility checks\\
Unassessable measured state & Native model information is insufficient to instantiate the compared sets\\
\bottomrule\end{tabularx}
\end{table}

A floating LP residual, including a nonpositive result within tolerance,
is not an equality certificate. A verified violating witness establishes
nonemptiness and strict inclusion even if a separate route or diagnostic
query remains unresolved. Eighteen internal assessments retained rationally
verified violations despite exhaustion of the SMT budget. They remain
certified states rather than additional unresolved assessments.

Witness checking preserves the complete net power query, exclusivity,
nonnegative directional flows, initial condition and energy dynamics. It
then checks every native energy, power, coupling and terminal row. A nonzero
exact residual violates an equality; its sign is not silently discarded.
Saved SMT assignments use exact fractional arithmetic. LP candidate flows
are converted from decimal strings to rationals; a further recorded attempt
limits denominators to $10^9$. Energy is recomputed with the same source
coefficients. Acceptance requires exact relaxed feasibility and a positive
reconstructed violation. Conversion does not repair an infeasible candidate
or introduce a feasibility tolerance.

Some external device instances also admit exact structural certificates
or constructive rational witnesses before the SMT fallback. In a lossless
constant efficiency device, cancellation preserves energy. Nonnegative
charge plus discharge coefficient sums then preserve the retained native
inequalities; equality rows require zero sums. A feasible idle trajectory
establishes nonemptiness in the implemented certificate. Constructive lossy
witnesses are accepted only after checking the complete mapped native system.

\subsection{Recorded settings and resource limits}
Table~\ref{tab:solver_settings} records the numerical settings and hardware used for the horizon study.

\begin{table}[htbp]
\caption{Software and numerical settings of the horizon study.}\label{tab:solver_settings}
\begin{tabularx}{\linewidth}{@{}p{.43\linewidth} X@{}}\toprule
Setting & Recorded value\\\midrule
Python / NumPy / SciPy & 3.12.14 / 2.3.5 / 1.16.2\\
HiGHS / Z3 & 1.8.0 / 4.15.3\\
SciPy LP method / presolve & \texttt{highs} / enabled\\
LP threads & 1 per call\\
Primal / dual feasibility tolerance & $10^{-7}$ / $10^{-7}$\\
Interior point optimality option & $10^{-8}$, when that backend applies\\
Positive floating residual threshold & $10^{-8}$ on the native row scale\\
Floating route comparison tolerance & $10^{-7}$ absolute\\
Z3 random seed / memory option & 0 / 768 MB\\
Shared SMT budget & 300 s per state\\
Shared state working budget & 1800 s from model construction\\
Reduced subproblem cap & 262144 planned subproblems per state\\
Full sign reference & Evaluated only at $T\leq8$\\
Exact rational acceptance tolerance & None\\
Processor & Intel Core i 9-13900HX; 24 cores,32 logical processors\\
Recorded physical memory & 16,907,100,160 bytes\\
Operating system & Windows 11, build 26200\\
\bottomrule\end{tabularx}
\end{table}

The serial order is SMT, reduced LP, then the full sign reference when
eligible. The SMT feasibility, violation and diagnostic queries share the
300 s allowance. The state deadline includes construction and preprocessing;
the 1800 s state allowance and 300 s SMT allowance are not additive.
Each LP uses the remaining state time, with a minimum passed time limit of
0.001 s after a positive remaining-time check. The reduced route is not
started if its predicted subproblem count exceeds the cap.

Escalation through the pressure horizons proceeds unless every eligible
state has only unresolved or budget limited routes and no exact certificate.
The extension beyond 96 intervals is entered after an exact certificate or a
completed reduced LP result at the preceding horizon. Two consecutive
extension horizons with neither stop further extension. Unentered design
slots remain separately accounted for. In the retained study, all 17 stated
horizons were reached.

The mapped device SMT fallback has a separate 300 s deadline starting after
attempted structural certification. Earlier native row checks used 30000 ms
per solver check, and the earlier affine row pilot used a coefficient sign
threshold of $10^{-12}$. These historical settings are not overwritten with
the horizon study settings. Exact classifications always retain the
applicable implementation and certificate basis.

\subsection{Timing design and noncompleted computations}
Original state wall times include construction, preprocessing and all
prescribed routes, not just time to the first certificate. Timings describe
the recorded machine and implementation. SMT performs an additional
diagnostic existence test, while the full sign reference performs an
additional constrained optimization. Their raw time ratios are not a
controlled solver speedup comparison.

Figure~\ref{fig:repeated_timing} shows the within input timing variation.
The repeated panel contains 36 fixed inputs at horizons 4,8,12,16,24 and 32.
For each horizon, all three efficiencies and both row forms are represented.
Using zero based efficiency and row indices $i,j$, duration is selected from
$(1/60,1/4)$~h by $(i+j)\bmod2$, and the before/after state by
$(i+j+T/4)\bmod2$. Source initialization is used throughout. Selection
preceded the observed timing results.

Each input receives one warmup and five measured repetitions under the
same options. The companion contains 108 input and route records and all
scheduled outcomes. Of 540 scheduled route repetitions,365 completed,
55 timed out and 120 full sign repetitions were skipped by the horizon rule.
A recorded zero for a skipped route is an availability placeholder, not a
measured time. Completed-only summaries are identified as such. Repetitions
measure timing variability of the same inputs and add no scientific states.


Markers give medians and bars the range of completed repetitions. R1 is reduced LP, R2 exact SMT and R3 the full sign reference. Timed out and skipped repetitions remain in the source records and are not plotted as completed times.

FIGINSERT1END



\subsection{Joint factor structure of the lossless control}
Figure~\ref{fig:lossless_factors} retains the joint classifications that are
not identifiable from horizon totals alone. Under the original protocol, the 312 lossless state slots
comprise 144 equality certificates, 104 unresolved assessments and 64
inapplicable assessments; no strict inclusion is certified in this control.
The marked factor combinations distinguish interval duration, net or gross
power rows, initialization and the before or after state. A dash denotes a
combination outside the prescribed long horizon panel, rather than another
assessed state. The complete lookup is supplied separately from the figure. The 104 originally unresolved lossless inputs are resolved by the later structural certificates described in Section~\ref{sec:sat_certificate_followup}.


N and G identify net and gross power rows. Source and boundary identify initialization rules. A dash denotes an unassessed combination outside the prescribed design. This figure retains the original protocol outcomes; later structural certificates resolve the originally unresolved inputs.

FIGINSERT2END



\section{Exact convex hull construction and numerical controls}\label{sec:hull}
The construction in Main Section~3.4 uses the net power formulation of
Section~\ref{sec:transitions}, specialized to two intervals. Initial energy
is $e=0.519$~MWh, duration is $h=0.25$~h, and the endpoint amplitude is
$a\in\{0,0.05,0.10\}$~MW. The nine cases cross these amplitudes with
$\eta_c\in\{0.65,0.90,1\}$. The remaining bounds, auxiliary drain,
discharging efficiency and native start/end limits are unchanged.

\subsection{Complete exclusive endpoints}
Write $\eta=\eta_c$, $U=0.57$~MWh and $\beta=50/3$~h$^{-1}$.
The first endpoint charges at amplitude $a$ in interval one; the second
discharges at the same amplitude. Their first energies are the values in
Main Section~3.4. Both endpoints charge exclusively in interval two.
Imposing the final charging boundary gives the intermediate calculation
\begin{align}
c_2^\nu+\beta\bigl(E_1^\nu+h(\eta c_2^\nu-P_{\mathrm{aux}})\bigr)
 &=\beta U,\\
(1+\beta h\eta)c_2^\nu
 &=\beta(U-E_1^\nu+hP_{\mathrm{aux}}),\qquad \nu\in\{+,-\}.
\label{eq:endpoint_construction}
\end{align}
The positive denominator determines $c_2^\nu$ uniquely. Reaching this
boundary alone is insufficient to establish endpoint feasibility. With
$E_0^\nu=e$, $k_\nu=c_2^\nu$ and
$E_2^\nu=E_1^\nu+h(\eta k_\nu-P_{\mathrm{aux}})$, the remaining
conditions can be written compactly in charging positive power coordinates:
\begin{equation}\begin{aligned}
0&\leq a\leq P,\qquad E_{\min}\leq E_1^\nu\leq U,\\
-\beta_d(E_s^\nu-E_{\min})&\leq u_1^\nu
\leq\beta(U-E_s^\nu),\quad s=0,1,\label{eq:endpoint_first}\end{aligned}\end{equation}
\begin{equation}\begin{aligned}
0&\leq k_\nu\leq P,\qquad E_{\min}\leq E_2^\nu\leq U,\\
-\beta_d(E_s^\nu-E_{\min})&\leq k_\nu
\leq\beta(U-E_s^\nu),\quad s=1,2.\label{eq:endpoint_second}
\end{aligned}\end{equation}
Here $u_1^+=a$ and $u_1^-=-a$; the manuscript net powers are their
negatives. These are substitutions into the native constraints, not
additional restrictions invented for the example. Every saved endpoint
satisfies them, the shared power bound and the energy recursion.

\subsection{Midpoint, reconstruction and residual}
Average the complete endpoint vectors, including stored energy.
Affine dynamics and native relaxed inequalities preserve this midpoint,
which belongs to the exact convex hull of the exclusive set. The first
flows are both $a/2$, while the second charge is
$(c_2^++c_2^-)/2$ and the second discharge is zero.
Exclusive reconstruction removes the equal first flows and retains the
second interval power. Thus
\begin{equation}
E_1^*-\bar E_1=E_2^*-\bar E_2
=\frac{h(1-\eta)a}{2}.\label{eq:hull_energy_shift}
\end{equation}
The midpoint final charging row is active, so its reconstructed excess is
\begin{equation}
\bar c_2+\beta E_2^*-\beta U
=\beta(E_2^*-\bar E_2)
=\frac{\beta h(1-\eta)a}{2}.\label{eq:hull_power_shift}
\end{equation}
The energy difference is in MWh and the power row excess is in MW.
With feasible endpoints, positive amplitude and charging losses give a
strictly positive excess. Reconstructed energy feasibility is checked
separately and holds for all nine stored constructions.

Table~\ref{tab:residual_controls} gives the energy shifts and native charging excesses for all nine controls.
Each midpoint is the equal weight combination of two feasible exclusive
trajectories. The table records the first interval energy increase
$\Delta E_1$ and the second interval charging constraint residual $\Delta P_2$.
All reconstructed energy trajectories satisfy their bounds.

\begin{table}[htbp]
\caption{Convexification controls.}\label{tab:residual_controls}

\begin{tabular}{@{}c c c c@{}}
\toprule
Charging efficiency & Amplitude (MW) & $\Delta E_1$ (kWh) & $\Delta P_2$ (kW)\\
\midrule
$0.65$ & $0$ & $0$ & $0$ \\
$0.65$ & $0.05$ & $2.1875$ & $875/24$ \\
$0.65$ & $0.1$ & $4.375$ & $875/12$ \\
$0.9$ & $0$ & $0$ & $0$ \\
$0.9$ & $0.05$ & $0.625$ & $125/12$ \\
$0.9$ & $0.1$ & $1.25$ & $125/6$ \\
$1$ & $0$ & $0$ & $0$ \\
$1$ & $0.05$ & $0$ & $0$ \\
$1$ & $0.1$ & $0$ & $0$ \\
\bottomrule
\end{tabular}
\end{table}

Four constructions have a positive residual; five have zero residual.
The zero cases arise from zero amplitude or unit charging efficiency.
They establish feasibility of the stored midpoint reconstruction,
not equality of the complete query sets.


\subsection{Linear objectives and complete query realizability}
For any linear objective on the complete variable vector, its value at a
convex combination is the same combination of endpoint values. Taking a
supremum over the exclusive set or its convex hull therefore gives the
same value whenever that value is well defined. With an attained optimum
this is equality of optimal values. The same statement applies to linear
objectives of the projected query.

This property does not require every projected hull point to lie in the
original query set. In the stored construction, the unique exclusive lift
of the midpoint query violates a native row. Uniqueness rules out another
reference realization with that query. It follows that equality of all
linear support values can coexist with different query sets. Agreement
for only eight selected objectives provides a still narrower comparison.

In the same complete trajectory coordinates, any convex outer approximation
containing the exclusive set also contains the two feasible endpoints and
their midpoint. It therefore retains this nonrealizable query. Additional
valid convex cuts in those coordinates cannot exclude the midpoint while
retaining every exclusive trajectory. This is a direct containment
consequence of the construction, with the same declared linear query;
it does not extend the statement to a changed query or reference model.


\section{Temporal queries and objective comparisons}
\label{si:temporal_queries}

\subsection{Model conditions and query projections}
\label{si:temporal_setup}
The comparisons specialize the model pair and reconstruction defined in the main Methods to the BOX formulation used for the analytical validity domains. There are two intervals of $h=0.1$~h, with $\eta_c=\eta_d=0.9$, unit retention, zero energy offset, energy bounds $[0,1]$~MWh and the shared power limit $c_t+d_t\leq1$~MW. Initial energy is fixed during each assessment, and terminal energy is free within its bounds. No additional coupling constraint is imposed.

The complete query is $p=(p_1,p_2)$, where $p_t=d_t-c_t$. Its scalar projections are $p_1$, $p_2$, $p_1+p_2$ and $p_1-p_2$. The corresponding equality domains are given in Table~\ref{tab:si_temporal_domains}. Here equality (EQ) means that the relaxed and exclusive models admit exactly the same values of the specified query; strict inclusion (STRICT) means that the relaxed image contains an additional value.

For the derivations, use dimensionless coordinates $e=E_0/E_s$, $x_t=p_t/P_s$ and $v_t=w_t/P_s$, with $E_s=1$~MWh, $P_s=1$~MW and $\tau=hP_s/E_s=1/10$. The simultaneous component satisfies $0\leq v_t\leq(1-|x_t|)/2$. Define $\eta=9/10$, $\kappa=1/\eta-\eta=19/90$ and $\alpha=(\eta+1/\eta)/2=181/180$. The energy reconstruction identity in the main Methods reduces to
\begin{equation}
\frac{E_t^*}{E_s}
=e-\tau\sum_{s=1}^{t}\left(\alpha x_s+\frac{\kappa}{2}|x_s|\right),
\qquad
E_t^*-E_t=\tau E_s\kappa\sum_{s=1}^{t}v_s.
\label{eq:si_temporal_reconstruction}
\end{equation}

\begin{table}[t]

\caption{Query preservation domains.}
\label{tab:si_temporal_domains}
\begin{tabularx}{\linewidth}{@{}l X@{}}
\toprule
Query & Initial stored energy preserving equality (MWh) \\
\midrule
$p_1$ & $[0,0.91]$ \\
$p_2$ & $[0,1]$ \\
$p_1+p_2$ & $[0,0.82]$ \\
$p_1-p_2$ & $[0,0.91]$ \\
$(p_1,p_2)$ & $[0,0.82]$ \\
\bottomrule
\end{tabularx}
\par\smallskip
\begin{minipage}{\linewidth}
The model conditions are specified in Section~\ref{si:temporal_setup}. Each scalar query is a projection of the complete two interval net power trajectory. All interval endpoints are included.
\end{minipage}
\end{table}



\subsection{Exact query domains}
\label{si:temporal_domains_proof}
For the first period, the exclusive and relaxed scalar images, in units of $P_s$, are respectively
\begin{equation}\begin{aligned}
\mathcal I_{1, X}(e)&=\left[-\min\left\{1,\frac{1-e}{\tau\eta}\right\},\ \min\left\{1,\frac{e\eta}{\tau}\right\}\right],\label{eq:si_p1X}\end{aligned}\end{equation}
\begin{equation}\begin{aligned}
\mathcal I_{1, R}(e)&=\left[-\min\left\{1,\frac{1-e+\tau\kappa/2}{\tau\alpha}\right\},\ \min\left\{1,\frac{e\eta}{\tau}\right\}\right].\label{eq:si_p1R}
\end{aligned}\end{equation}
At fixed charging net power, the smallest relaxed energy is obtained at the largest permitted simultaneous component. This gives the lower endpoint in Eq.~\eqref{eq:si_p1R}; pure discharging gives the common upper endpoint. The charging endpoints agree up to $e=1-\tau\eta=0.91$ and differ strictly above it.

Both models have the same second period image,
\begin{equation}
\mathcal I_2(e)=\left[-1,\ \min\left\{1,\frac{e\eta}{\tau}+\eta^2\right\}\right].
\label{eq:si_p2}
\end{equation}
The extreme first period energies, $\max\{0,e-\tau/\eta\}$ and $\min\{1,e+\tau\eta\}$, are attainable exclusively. The minimum permits full charging in period two, while the maximum determines the largest discharge. Simultaneous operation cannot extend either required energy extreme. Within each power sign region the feasible scalar image is an interval containing the idle query, so these endpoints determine the complete scalar images.

For the complete query, removing simultaneous operation raises energy and reduces the shared power use. Lower energy bounds therefore remain satisfied. If $e\leq1-2\tau\eta=0.82$, reconstructed energy cannot exceed $e+2\tau\eta\leq1$, proving complete query equality on this entire closed interval. For $0.82<e\leq0.91$, define
\begin{equation}
C_X(e)=\frac{1-e}{\tau\eta},\qquad
C_R(e)=\frac{1-e+\tau\kappa}{\tau\alpha}.
\label{eq:si_temporal_cuts}
\end{equation}
Then $C_R>C_X$. The query $x_1=x_2=-C_R/2$ is relaxed feasible when $v_t=(1-|x_t|)/2$: its energies are $(e+1)/2$ and $1$. Its exclusive terminal energy is $e+\tau\eta C_R>1$. For $e>0.91$, Eqs.~\eqref{eq:si_p1X}--\eqref{eq:si_p1R} already give a first period query available only to the relaxation; setting the second period powers to zero extends it to a full trajectory. Thus the complete equality domain is exactly $[0,0.82]$~MWh, consistent with the complete-query domain in Table~\ref{tab:boundary_domains}.

For the sum query, the common upper endpoint is $\min\{2,e\eta/\tau\}$. Its relaxed lower endpoint is $-\min\{2,C_R(e)\}$, while the exclusive lower endpoint is $-C_\Sigma(e)$, where
\begin{equation}
C_\Sigma(e)=
\begin{cases}
2,&0\leq e\leq 0.82,\\
(1-e)/(\tau\eta),&0.82<e\leq 0.91,\\
1-\eta^2+\eta(1-e)/\tau,&0.91<e\leq 1.
\end{cases}
\label{eq:si_sum_endpoint}
\end{equation}
The last branch uses a first period discharge of $\eta(e-0.91)/\tau$ followed by unit charging. Comparing the four exclusive mode sequences shows that this gives more total net charging than two charging periods when $e>0.91$. A final discharge cannot improve the minimum sum, and two discharges give a nonnegative sum. The sum images therefore agree exactly for $e\leq0.82$.

For the difference query, the common upper endpoint is $1+\min\{1,e\eta/\tau\}$. For $e\leq0.91$, the common lower endpoint is $-1-\min\{1,\eta(e+\tau\eta)/\tau\}$, attained by full charging followed by the largest feasible discharge. Above $0.91$, it equals the respective first period lower endpoint in Eqs.~\eqref{eq:si_p1X}--\eqref{eq:si_p1R}, minus one. These facts establish the remaining equality domain in Table~\ref{tab:si_temporal_domains} without interpolation between evaluated states.

\subsection{Directional objective comparisons}
\label{si:temporal_objectives}
On $e\in[0.82,0.91]$, the exclusive query set in dimensionless coordinates is the square $[-1,1]^2$ intersected with $x_1+x_2\geq-C_X(e)$. The relaxed query set has the corresponding cut $-C_R(e)$. These cuts affect only the corner with two charging periods because both constants lie in $[1,2]$. The first period upper energy bound and all lower energy bounds are satisfied throughout this square on the stated slice. For the relaxation, maximal simultaneous components give the minimum terminal energy $e-\tau\alpha(x_1+x_2)-\tau\kappa$, which proves the relaxed projection formula.

For the objective $\max[-wp_1-(1-w)p_2]$, $w\in[0,1]$, the optimal value for a cut constant $C$ is
\begin{equation}
P_s\bigl[\max\{w,1-w\}+\min\{w,1-w\}(C-1)\bigr].
\end{equation}
The relaxed optimum minus the exclusive optimum is consequently
\begin{equation}
\Delta(E_0, w)=\min\{w,1-w\}\, K(E_0-E_c),\quad
K=\frac{1900}{1629}\,\mathrm{MW\, MWh^{-1}},\quad E_c=0.82~\mathrm{MWh}.
\label{eq:si_temporal_gap}
\end{equation}
Equation~\eqref{eq:si_temporal_gap} applies only to $E_0\in[0.82,0.91]$~MWh. At $w=0$ or $1$ the gap is zero throughout this slice, although complete query equality fails for $E_0>0.82$~MWh. At $E_0=0.90$~MWh and $w=1/2$, it is $76/1629$~MW, approximately $46.6544$~kW.

The directional comparisons use eight vectors $u$ with $\|u\|_1=1$ and plot the gaps between the maxima of $u^\mathsf{T}p$. Their order follows objective angle. They are selected objective comparisons, rather than scores of model quality. Zero gaps in a finite set of directions do not establish complete query equality. Even equality of every linear support function identifies convex hulls unless further set structure is established.

Table~\ref{tab:directions} lists all eight coefficient vectors and the 24 displayed gaps, including every zero. Figure~\ref{fig:si_temporal_response} separates the objective response from a selected trajectory witness.
\begingroup
\begin{longtable}{@{}p{0.12\linewidth}p{0.09\linewidth}p{0.09\linewidth}p{0.19\linewidth}p{0.19\linewidth}p{0.19\linewidth}@{}}
\caption{Coefficients and objective gaps for all displayed directional objectives.}\label{tab:directions}\\
\toprule Direction & $u_1$ & $u_2$ & 0.82 MWh (kW) & 0.90 MWh (kW) & 1.00 MWh (kW)\\
\midrule\bottomrule
D1 & 1 & 0 & 0 & 0 & 0\\
D2 & 0.5 & 0.5 & 0 & 0 & 0\\
D3 & 0 & 1 & 0 & 0 & 0\\
D4 & -0.5 & 0.5 & 0 & 0 & 52.486188\\
D5 & -1 & 0 & 0 & 0 & 104.972376\\
D6 & -0.5 & -0.5 & 0 & 46.654389 & 9.972376\\
D7 & 0 & -1 & 0 & 0 & 0\\
D8 & 0.5 & -0.5 & 0 & 0 & 0\\
\end{longtable}
 The gap is the relaxed maximum minus the exclusive maximum of the same linear objective. All 24 entries, including zeros, are retained. Values are rounded for display; the companion data retain exact fractional gaps and both optimal values. The zero gap counts are 8,7 and 5 for the three energy states, respectively. Lines in the radar display join discrete objectives only.
\endgroup



Panel (a) shows the linear objective gap over the stated initial energy and weight slice. Panel (b) shows a selected feasible relaxed query with an 8 kWh upper energy violation after reconstruction. The trajectory in panel (b) is not asserted to optimize panel (a).

FIGINSERT3END


\subsection{A native energy violation witness}
\label{si:temporal_witness}
At $E_0=0.90$~MWh, choose $c_1=c_2=0.8$~MW and $d_1=d_2=0.2$~MW. The shared power limit is satisfied, and the relaxed energies are $2137/2250$ and $2249/2250$~MWh. Both net powers are $-0.6$~MW. Their unique exclusive reconstruction has $c_1^*=c_2^*=0.6$~MW and $d_1^*=d_2^*=0$, with energies $477/500$ and $126/125$~MWh. The relaxed trajectory is feasible, whereas reconstruction exceeds the upper energy bound by $1/125$~MWh, or $8$~kWh (Fig.~\ref{fig:si_temporal_response}b).


This is a selected feasible-query witness, not necessarily an optimizer of the displayed objective. Table~\ref{tab:si_temporal_values} separately reports the largest reconstructed upper energy residual over the relaxed feasible set. The failure here concerns an energy bound; it is distinct from the native power constraint violation with admissible reconstructed energy reported in the main Results.

\begin{table}[t]

\caption{Representative query comparisons.}
\label{tab:si_temporal_values}


\begin{tabularx}{\linewidth}{@{}l l l l X X@{}}
\toprule
$E_0$ (MWh) & $p_1$ & $p_2$ & Full query & Objective gap (kW) & Energy excess (kWh) \\
\midrule
0.82 & EQ & EQ & EQ & 0.000 & 0.000 \\
0.90 & EQ & EQ & STRICT & 46.654 & 8.398 \\
0.91 & EQ & EQ & STRICT & 52.486 & 9.448 \\
1.00 & STRICT & EQ & STRICT & 9.972 & 18.895 \\
\bottomrule
\end{tabularx}
\par\smallskip
\begin{minipage}{\linewidth}
The objective is $\max[-(p_1+p_2)/2]$; its gap is the relaxed optimum minus the exclusive optimum. Energy excess is the nonnegative maximum reconstructed upper energy residual over the relaxed feasible set. At $E_0=0.90$~MWh this worst excess differs from the selected witness's $8$~kWh violation. Exact rational values are retained in the source data.
\end{minipage}
\end{table}



\subsection{Extension to finite horizons}
\label{si:temporal_horizon}
Retain the same coefficients, free terminal energy and $E_0=0.90$~MWh, but allow any finite horizon $T$. Every first period exclusive power in $[-1,1]$~MW gives energy in $[71/90,99/100]$~MWh. For any later target period, discharge at $1$~MW in the first period and remain idle until the target. Energy before the target is $71/90$~MWh. Applying any target power in $[-1,1]$~MW then gives energy in $[61/90,791/900]$~MWh. Remaining idle preserves feasibility. Thus every individual period has the same scalar image $[-1,1]$~MW under both models.

For every finite $T\geq2$, extend the preceding two period relaxed witness by idle intervals. Its energy remains $2249/2250$~MWh after period two, while its unique exclusive reconstruction remains at $126/125$~MWh. Hence the full query sets differ for every such horizon, despite equality of each individual period image. At $T=1$, the full query is scalar and equality holds. The argument depends on the stated retention, offset and terminal conditions; it does not extend unchanged to additional couplings or different dynamics.

\subsection{Reproducibility and interpretation}
The source data retain exact rational objective values, scalar interval endpoints, feasible relaxed vertices and reconstructed energy residuals. The displayed classification matrix comprises thirteen prescribed energy states and five queries. The directional profiles use eight fixed objective vectors, with three representative states shown. These are deterministic comparisons, rather than independent devices or operating frequencies. A separate reconstruction in net power and simultaneous-component coordinates verifies the saved values. Analytical boundaries and the finite horizon statement follow from the preceding proofs, rather than from the number of verification points or the rendering grid.


\section{Source specific device mappings}\label{sec:external}
\subsection{Common coordinates and source scope}
Each executable mapping defines a fixed device instance, a relaxed set and
its exclusive reference. Initial stored energy is the conditioning state;
the published or documented ratings, efficiencies, duration, horizon and
terminal rule remain fixed within that instance. The query retains the
complete discharge positive power trajectory. Independent price, network,
investment and policy decisions are excluded when they cannot be represented
by the documented device constraints. A source's physical experiment does
not change the evidence level of this separate mathematical device assessment.

All instances have nonnegative directional powers, separate ratings and
bounded stored energy. In their documented source units the common dynamics are
\begin{equation}
E_t=E_{t-1}+h(\eta_cc_t-d_t/\eta_d),\qquad
0\leq c_t\leq P_c,\quad 0\leq d_t\leq P_d .\label{eq:external_common}
\end{equation}
Retention is one and the selected device equations have zero auxiliary
offset. Source units are retained: kW with kWh, or MW with MWh.
The source of each numerical parameter and the omitted components are
specified in Table~\ref{tab:source_locators}. Table~\ref{tab:external_parameters}
gives every executed instance rather than only representative parameters.
Instance identifiers connect these tables to the accompanying state records.

The box family retains the dynamics, separate power ratings and energy
bounds. The power convex hull family additionally retains
\begin{equation}
\frac{c_t}{P_c}+\frac{d_t}{P_d}\leq 1 .\label{eq:external_hull}
\end{equation}
The headroom family also retains the beginning of interval bounds
\begin{equation}
h\eta_cc_t+E_{t-1}\leq E_{\max},\qquad
\frac{h}{\eta_d}d_t-E_{t-1}\leq-E_{\min} .\label{eq:external_headroom}
\end{equation}
The valid inequality family adds the source reachable power limits and
multiperiod linear rows. Their coefficients depend on the documented
ratings, energy range, efficiency, initial state and remaining periods.
They are retained as linear rows at each fixed assessed state. Conic
extensions are excluded. The accompanying model builder records the exact
coefficient construction and fixes a directional flow to zero when its
available bound is zero, rather than introducing division by a small positive
number. Every mapped reference adds charge and discharge exclusivity to the
same retained native rows.

\begin{landscape}
\begingroup
\begin{longtable}{@{}p{0.16\linewidth}p{0.27\linewidth}p{0.29\linewidth}p{0.21\linewidth}@{}}
\caption{Source equations, numerical provenance and components outside the mapped device.}\label{tab:source_locators}\\
\toprule Source & Retained equations or source locator & Numerical parameters and transformations & Components outside the comparison\\
\midrule\bottomrule
Pozo~\cite{pozo2023convex} & Energy headroom formulation; implemented with the documented Morales device, Eqs.~(1)--(4) and the corresponding headroom rows. & First storage device, R000, in the Morales numerical example and accompanying author code. Numerical provenance is shared with the two Morales instances. & Ramps, reserves, network, objective and fleet operation.\\
Morales~\cite{morales2025linear} & Eqs.~(1)--(4); Theorem~1, Eqs.~(8)--(11). Natural relaxation and linear cuts are separate instances. & First storage device R000 in the source numerical example and author code. Only rows of the selected linear formulation are retained. & Network, objective, fleet operation and conic extensions.\\
Prat et al.~\cite{prat2025sufficient} & Eqs.~(1)--(5), Section~3.3 and Table~1, pp.~8--9. & First illustrative instance and three tabulated modes: lossless, lossy and slow. Selected leakage is zero. & Prices, objective dependent optimality and dispatch policy.\\
Nazir and Almassalkhi~\cite{nazir2021guaranteeing} & Constant efficiency dynamics and limits of the relaxed baseline in Section~V. & Published numerical baseline with initial energy assessed over $[0,60]$~kWh. Charging positive source power is negated for the query. & Robust inner dispatch approximation and its recovery guarantee.\\
Elgersma et al.~\cite{elgersma2026tight} & Basic operation, Eq.~(1), and time indexed operation, Eq.~(4). Retained 2024 preprint: Section~IV-A, p.~6; Table~IV, operation column, p.~9. & Two hourly intervals. Operation ratings are evaluated from source expressions. Eliminating the relaxed mode variable gives the normalized sum power bound exactly. & Network, reserves and investment; no extension of a single period hull theorem to the complete horizon.\\
Elsaadany and Almassalkhi~\cite{elsaadany2023battery} & Eqs.~(4a)--(4e) and (6b), Table~I and Section~V-C. & Three hour signal at five minute resolution: 36 intervals. The inclusive source index denotes these 36 intervals. & Objective dependent results and the full robust optimization procedure.\\
Chancelier et al.~\cite{chancelier2026optimal} & Section~3.1, Eqs.~(11b)--(11e) and (12b)--(12e); Section~5.2 and Fig.~5. & Published modes 0.4 and 1. Source coefficients 0.95 and 1.05 give efficiencies 0.95 and $20/21$. Delivery energy is divided by duration and negated for net power. & Cash, stochastic nonanticipativity, expected utility and trading policies.\\
Tziovani et al., successive convexification~\cite{tziovani2023successive} & Author manuscript, pp.~3--4, Eqs.~(1), (2) and (11)--(16); Tables~II, III and V; extended device, pp.~10--11. & Six devices at 24 and five intervals, plus the extended 480 interval case at $h=0.05$~h. The latter uses translated usable energy and efficiency 0.96. & Network, unit commitment, objectives, quadratic losses and trust region iterations.\\
Tziovani et al., piecewise linear comparison~\cite{tziovani2023investigating} & Eqs.~(1b), (9a), (9b) and (12); exclusive constraint (8b); Tables~II and III, Section~V-A. & Six source devices over 24 hourly intervals. Native dynamics and box limits are retained; complementarity is removed. & Network, unit commitment and objective. Other loss variants are not counted as new mappings.\\
Tziovani et al., buildings~\cite{tziovani2025energy} & Eqs.~(2), (3a), (3b) and (4); Sections~V-A and VI-C. & Daily device: 5~kWh, efficiency 0.96. Experimental approximation: 9.8~kWh, efficiency 0.91. The latter retains the 5~kW rating carried from Section~V. SOC is converted using capacity; terminal floors are fixed. & Network, monetary objective, degradation allocation, controller softening and nonlinear efficiency.\\
\end{longtable}
 Locators identify the primary source versions used for mapping. Pagination can differ in a later version of record, including the published Elgersma article. The Pozo and Morales instances share numerical provenance. Identifiers in Table~\ref{tab:external_parameters} distinguish executed formulations and parameter settings, rather than independent device datasets.
\endgroup
\end{landscape}

\begin{landscape}
\begingroup
\setlength{\LTpost}{3pt}
\begin{longtable}{@{}p{0.045\linewidth}p{0.28\linewidth}p{0.09\linewidth}p{0.09\linewidth}p{0.1\linewidth}p{0.09\linewidth}p{0.04\linewidth}p{0.07\linewidth}p{0.09\linewidth}@{}}
\caption{Parameters and outcomes of the 34 mapped device instances.}\label{tab:external_parameters}\\
\toprule ID & Operating instance & Power / energy units & Energy bounds & $P_c; P_d$ & $\eta_c; \eta_d$ & $T$ & $h$ (h) & EQ; STRICT\\
\midrule\bottomrule
1 & Pozo: energy headroom & kW/kWh & $[30,60]$ & 20; 20 & 0.9; 0.95 & 24 & 1 & 0; 111\\
2 & Morales: natural relaxation & kW/kWh & $[30,60]$ & 20; 20 & 0.9; 0.95 & 24 & 1 & 0; 111\\
3 & Morales: linear cuts & kW/kWh & $[30,60]$ & 20; 20 & 0.9; 0.95 & 24 & 1 & 0; 111\\
4 & Prat: first example & MW/MWh & $[0,1]$ & 1.5; 1.5 & 0.9; 0.9 & 24 & 1 & 0; 111\\
5 & Prat: lossless mode & MW/MWh & $[0,1]$ & 0.6; 0.6 & 1; 1 & 24 & 1 & 111; 0\\
6 & Prat: lossy mode & MW/MWh & $[0,1]$ & 0.6; 0.6 & 0.9; 0.9 & 24 & 1 & 0; 111\\
7 & Prat: slow mode & MW/MWh & $[0,1]$ & 0.2; 0.2 & 0.9; 0.9 & 24 & 1 & 0; 111\\
8 & Nazir: relaxed baseline & kW/kWh & $[0,60]$ & 15; 15 & 0.95; 0.95 & 24 & 1 & 0; 111\\
9 & Elgersma: basic operation & MW/MWh & $[5,13]$ & $80/9$; 7.2 & 0.9; 0.9 & 2 & 1 & 0; 111\\
10 & Elgersma: time indexed operation & MW/MWh & $[5,13]$ & $80/9$; 7.2 & 0.9; 0.9 & 2 & 1 & 2; 107\\
11 & Elsaadany: device relaxation & kW/kWh & $[0,13]$ & 5; 5 & 0.95; 0.95 & 36 & $1/12$ & 0; 111\\
12 & Chancelier: mode 0.4 & MW/MWh & $[0,1]$ & 0.4; 0.4 & 0.95; $20/21$ & 24 & 1 & 0; 111\\
13 & Chancelier: mode 1 & MW/MWh & $[0,1]$ & 1; 1 & 0.95; $20/21$ & 24 & 1 & 0; 111\\
14 & Tziovani: device 1, 24 intervals & MW/MWh & $[1,4]$ & 5; 5 & 0.89; 0.89 & 24 & 1 & 0; 111\\
15 & Tziovani: device 2, 24 intervals & MW/MWh & $[3,6.5]$ & 5.5; 5.5 & 0.91; 0.91 & 24 & 1 & 0; 111\\
16 & Tziovani: device 3, 24 intervals & MW/MWh & $[0.5,1.5]$ & 1.5; 1.5 & 0.88; 0.88 & 24 & 1 & 0; 111\\
17 & Tziovani: device 4, 24 intervals & MW/MWh & $[0.5,1]$ & 0.5; 0.5 & 0.92; 0.92 & 24 & 1 & 0; 111\\
18 & Tziovani: device 5, 24 intervals & MW/MWh & $[0.5,0.7]$ & 0.5; 0.5 & 0.89; 0.89 & 24 & 1 & 0; 111\\
19 & Tziovani: device 6, 24 intervals & MW/MWh & $[0.5,0.7]$ & 0.5; 0.5 & 0.91; 0.91 & 24 & 1 & 0; 111\\
20 & Tziovani: device 1, 5 intervals & MW/MWh & $[1,4]$ & 5; 5 & 0.89; 0.89 & 5 & 1 & 0; 111\\
21 & Tziovani: device 2, 5 intervals & MW/MWh & $[3,6.5]$ & 5.5; 5.5 & 0.91; 0.91 & 5 & 1 & 0; 111\\
22 & Tziovani: device 3, 5 intervals & MW/MWh & $[0.5,1.5]$ & 1.5; 1.5 & 0.88; 0.88 & 5 & 1 & 0; 111\\
23 & Tziovani: device 4, 5 intervals & MW/MWh & $[0.5,1]$ & 0.5; 0.5 & 0.92; 0.92 & 5 & 1 & 0; 111\\
24 & Tziovani: device 5, 5 intervals & MW/MWh & $[0.5,0.7]$ & 0.5; 0.5 & 0.89; 0.89 & 5 & 1 & 0; 111\\
25 & Tziovani: device 6, 5 intervals & MW/MWh & $[0.5,0.7]$ & 0.5; 0.5 & 0.91; 0.91 & 5 & 1 & 0; 111\\
26 & Tziovani: extended device & MW/MWh & $[0,1.2]$ & 2.4; 2.4 & 0.96; 0.96 & 480 & 0.05 & 0; 111\\
27 & Tziovani: box device 1 & MW/MWh & $[1,4]$ & 5; 5 & 0.89; 0.89 & 24 & 1 & 0; 111\\
28 & Tziovani: box device 2 & MW/MWh & $[3,6.5]$ & 5.5; 5.5 & 0.91; 0.91 & 24 & 1 & 0; 111\\
29 & Tziovani: box device 3 & MW/MWh & $[0.5,1.5]$ & 1.5; 1.5 & 0.88; 0.88 & 24 & 1 & 0; 111\\
30 & Tziovani: box device 4 & MW/MWh & $[0.5,1]$ & 0.5; 0.5 & 0.92; 0.92 & 24 & 1 & 0; 111\\
31 & Tziovani: box device 5 & MW/MWh & $[0.5,0.7]$ & 0.5; 0.5 & 0.89; 0.89 & 24 & 1 & 0; 111\\
32 & Tziovani: box device 6 & MW/MWh & $[0.5,0.7]$ & 0.5; 0.5 & 0.91; 0.91 & 24 & 1 & 0; 111\\
33 & Tziovani: daily building device & kW/kWh & $[0.75,4.75]$ & 5; 5 & 0.96; 0.96 & 24 & 1 & 0; 111\\
34 & Tziovani: experimental approximation & kW/kWh & $[3.43,8.33]$ & 5; 5 & 0.91; 0.91 & 480 & 0.05 & 0; 111\\
\end{longtable}
 Power and efficiency pairs list charging then discharging. All instances use unit retention and zero offset; initial energy varies within the listed bounds. Fixed terminal floors are 1.25~kWh (ID33) and 4.41~kWh (ID34), independently of assessed initial energy.
\endgroup
\end{landscape}


\subsection{Transformations and interpretation by source}
\paragraph{Pozo and Morales.}
The Pozo headroom formulation~\cite{pozo2023convex} uses the documented
numerical device also used for the Morales example. The two Morales
instances~\cite{morales2025linear} distinguish the natural linear
relaxation from the strengthened linear inequalities. Shared parameters
are not independent evidence from different devices. Source equation
locators that mention both variants do not mean that every row is active
in the natural relaxation. The retained model family determines the rows.

\paragraph{Prat and coauthors.}
Four source parameter settings cover the first illustrative case and the
three Table~1 modes~\cite{prat2025sufficient}. The selected leakage is zero.
No normalized sum power inequality is added to the box adapter. The lossless
mode has 111 equality states; the other three settings have 111 strict
inclusion states each. These are parameter instances of one source.

\paragraph{Nazir and Almassalkhi.}
The evaluated instance is the constant efficiency relaxed baseline from
the source numerical section~\cite{nazir2021guaranteeing}. It is not the
robust inner dispatch envelope developed in that study. The source uses
charging positive net power, so the declared power query is its negative.
The bijective sign change preserves the question of equality of query sets.

\paragraph{Elgersma and coauthors.}
The basic operation and time indexed operation instances retain the source
operation parameters~\cite{elgersma2026tight}. Eliminating the relaxed mode
variable is exact because, for positive ratings,
\begin{equation}
\delta_t\in[c_t/P_c,\;1-d_t/P_d]
\quad\Longleftrightarrow\quad c_t/P_c+d_t/P_d\leq 1 .\label{eq:mode_elimination}
\end{equation}
The time indexed model also retains the energy headroom rows. Operation
ratings are evaluated from the published expressions, not rounded before
the assessment. The retained source version is documented in the mapping
locators. Network, reserves and investment are excluded. A source theorem
about a single period hull is not attributed to the full multiperiod query.

\paragraph{Elsaadany and Almassalkhi.}
The selected device relaxation retains the published sum power inequality,
as well as energy dynamics and bounds~\cite{elsaadany2023battery}.
The three hour signal with five minute intervals has 36 intervals. This
resolves the source indexing convention without adding a 37th interval.
The source robust optimization and objective conditioned guarantees are
outside the mapped device comparison.

\paragraph{Chancelier and coauthors.}
Two documented operating modes retain deterministic device paths from the
stochastic trading formulation~\cite{chancelier2026optimal}. Charging
positive delivery energy is divided by the published duration and negated
to obtain average discharge positive power. A discharged energy coefficient
of 1.05 corresponds to directional efficiency $20/21$, rather than 1.05.
Cash, utility and nonanticipative policies are excluded. The two modes
come from the source examples; the evaluated initial energies are not
reported field observations.

\paragraph{Tziovani sources.}
The earlier piecewise linear comparison contributes six box device instances
with its complementarity constraint removed~\cite{tziovani2023investigating}.
The successive convexification study contributes the documented 24 interval,
five interval and extended device cases~\cite{tziovani2023successive}.
Its linear loss coordinates map to directional powers through
\begin{equation}
P_t=d_t-c_t,\qquad
\ell_t=(1/\eta_d-1)d_t+(1-\eta_c)c_t .\label{eq:loss_coordinates}
\end{equation}
The retained loss envelope corresponds to nonnegative flows and the
normalized sum power bound. The extended device translates the energy
origin to the lower usable energy level. Its relative zero does not mean
physical zero SOC. The documented efficiency 0.96 is retained, rather than
treating a rounded displayed loss coefficient as a different exact parameter.

The building study contributes a daily model and a constant efficiency
approximation of its experimental device~\cite{tziovani2025energy}.
SOC fraction is multiplied by capacity to obtain energy. The final energy
lower bound is fixed independently of the initial energy being assessed.
The two native terminal rows are
\begin{equation}
E_T\geq 1.25~\mathrm{kWh}\quad\text{(daily device)},\qquad
E_T\geq 4.41~\mathrm{kWh}\quad\text{(experimental approximation)}.
\label{eq:terminal_floors}
\end{equation}
These bounds are fixed as initial energy varies. The executable terminal
adapter adds this row to the common device builder before verification.
The mapping retains that terminal floor and excludes degradation allocation,
controller softening and nonlinear efficiency. It does not certify the
measured device's complete native query set.

\subsection{Initial energy sampling and nested outcome units}
The initial grid uses 101 uniformly spaced inclusive states. With
$q=(E_0-E_{\min})/(E_{\max}-E_{\min})$, these are $q=i/100$ for
$i=0,\ldots,100$. For an instance with one common grid classification,
the ten additional fractions are
\begin{equation}
q\in\left\{\frac18,\frac58,\frac38,\frac78,\frac1{16},\frac9{16},
\frac5{16},\frac{13}{16},\frac3{16},\frac{11}{16}\right\}.
\label{eq:external_extra_states}
\end{equation}
For the time indexed Elgersma instance, the eight boundary refinements are
$q=1/(100\,2^k)$ and $q=1-1/(100\,2^k)$, for $k=1,2,3,4$.
The original endpoints are already in the inclusive grid. Repeated complete
model states are identified by the fixed parameters and initial energy
rather than counted again as new device instances.
The time indexed Elgersma panel has 109 retained states; the other 33
instances have 111 each. This gives 3772 state-instance assessments.
The 113 equality cases consist of 111 Prat lossless states and the two
Elgersma endpoints. The 3659 strict cases have retained exact witnesses.
Full state locations, certificate bases and preserved source scope are in
the companion data. The grid is a deterministic design, not a probability
distribution over storage operation.

\subsection{Elgersma endpoint proof and feasible updates}
For the mapped time indexed operation device, $h=1$~h,
$\eta_c=\eta_d=0.9$, the energy range is $[5,13]$~MWh,
$P_c=80/9$~MW and $P_d=36/5$~MW. At the lower initial endpoint,
first interval discharge headroom forces zero discharge. At the upper
endpoint, charge headroom forces zero charge. Removing simultaneous flows
therefore leaves the first energy unchanged. In interval two the exclusive
directional powers are no larger than the relaxed ones. The unchanged
preceding energy and retained headroom rows preserve the bounds. Both
endpoints have nonempty feasible sets and preserve the complete query.

For an interior energy $e$, let $L=5$ and $U=13$ in MWh. The saved
interior witness chooses positive discharge $D$ satisfying
\begin{equation}
D<\tfrac 14\min\{(e-L)\eta_d/h,\;(U-e)\eta_d/h,\;
P_d,\;P_c\eta_c\eta_d\},\qquad C=D/(\eta_c\eta_d).
\label{eq:elgersma_interior}
\end{equation}
Charging at $C$ and discharging at $D$ leaves the first relaxed energy
equal to $e$. Exclusive reconstruction increases it by
$hD(1/\eta_d-\eta_c)>0$. A second exclusive charge of
$(U-e)/(h\eta_c)$ reaches $U$ in the relaxed trajectory and exceeds it
after reconstruction. These flows satisfy the retained rating and
headroom conditions under the stated choice of $D$. This is the study's
query characterization for the mapped device, not a theorem attributed
to the source publication.

Table~\ref{tab:external_transitions} lists the four saved feasible updates.

\begin{table}[htbp]
\caption{Certified feasible updates at the mapped device endpoints.}\label{tab:external_transitions}
\begin{tabular}{@{}rrrrl@{}}\toprule
$E_0$ (MWh)&Updated (MWh)&Charge (MW)&Discharge (MW)&Query relation\\\midrule
5&5.005&$1/180$&0&EQ to STRICT\\
5.005&5&0&$9/2000$&STRICT to EQ\\
13&12.995&0&$9/2000$&EQ to STRICT\\
12.995&13&$1/180$&0&STRICT to EQ\\
\bottomrule\end{tabular}
\end{table}

Each action lasts one hour. The complete two interval query and all device
coefficients remain fixed between the corresponding assessments.
The recorded 0.005 MWh boundary brackets are discrete endpoint/interior
checks; the endpoint locations follow from the analytical argument.

Table~\ref{tab:residual_family_witnesses} separates energy and power failures in the retained witnesses.
\subsection{Witness categories by formulation}
The classification below links each mapped constraint family to the
failure categories of its retained witnesses. Counts refer to model states,
which remain nested within source formulations and operating instances.

\begin{table}[htbp]
\caption{Witness categories by formulation.}\label{tab:residual_family_witnesses}
\begin{tabular}{@{}l r r r@{}}
\toprule
Constraint family & Energy only & Power only & Energy and power\\
\midrule
Energy headroom & 0 & 0 & 218 \\
Power convex hull & 1776 & 0 & 0 \\
Valid inequalities & 0 & 0 & 111 \\
Box limits & 1554 & 0 & 0 \\
\bottomrule
\end{tabular}
\par\smallskip
\begin{minipage}{\linewidth}
Power categories refer to the retained device power or energy headroom
inequalities. Their residuals keep the units of the source row and do not
represent a common delivered power overload.
\end{minipage}
\end{table}

The 3659 model states represented by these counts are certified cases of strict
query set inclusion. Energy headroom and valid inequality cases retained
combined energy and power violations; box and power convex hull cases
retained energy violations alone. No retained external witness exhibited
power failure with an admissible stored energy trajectory. These categories
refer to the frozen witnesses rather than all possible witnesses or an
operational failure probability. The 113 equality cases are not included
in the failure counts.



Within the box and normalized power rating families, removing the common
charge and discharge component cannot violate a retained rating. In the
headroom family, exclusive next-interval energy feasibility implies the
corresponding charging or discharging headroom row. The source valid
inequalities likewise retain their validity for an exclusive device
trajectory satisfying the native energy and rating constraints. This gives
a structural explanation for the absence of energy feasible power failures
in these particular mappings. The state dependent taper constraints used in
Section~\ref{sec:transitions} impose a different relationship between power
and stored energy and are not covered by that argument.


\section{Simulated and measured operating states}\label{sec:operating_states}
\subsection{Simulation driven by public electricity prices}
The source driven simulation uses day ahead prices from the Danish DK1 bidding zone for
January 2025 and a documented device with energy limits 30 to 60 kWh,
charge and discharge ratings 20 kW, efficiencies 0.9 and 0.95, unit retention
and one hour intervals. The device retains the headroom formulation
described in Section~\ref{sec:external}. Its declared query horizon is four
intervals. The first 168 price observations define a right continuous
empirical cumulative distribution. Calibration observations are excluded
from the assessed windows.

Starting at 55 kWh, a bounded tracking rule maps the current price quantile
to a target energy between 30 and 60 kWh. At each interval it charges or
discharges exclusively toward that target, subject to the directional
20 kW rating and the same device dynamics. A higher price quantile sets a
higher energy target. This rule was selected for state coverage and does
not claim economic optimality. The original policy, input identity and
exact trajectory values are retained with the computational records.

The windows start on 8,16 and 25 January 2025 and each contain 168 hourly
actions plus their final energy timestamp. Their 507 timestamp records map
to 86 distinct complete model states. Each complete state has a saved strict
inclusion certificate. Repeated visits are not additional model states.
Public market inputs driving a simulated device remain separate from
measurements of a physical battery.

\subsection{Measurement fields and selection of windows}
Figure~\ref{fig:measured_coverage} displays the measured power and SOC coverage.
The measured windows originate from ACEN's New England BESS public
operating data and the Electric Power Research Institute (EPRI) PERIOD data resource. Source references and
access conditions are provided in the accompanying data guide. ACEN records
contain dashboard alternating current (AC) power and battery management system (BMS) SOC for two device
units. EPRI records contain direct current (DC) battery power, battery SOC and
the reported SOC limits. Source power scaling, measurement side and time
columns are preserved. A nominal capacity times BMS SOC is an energy
proxy, not an independent measurement of thermodynamic stored energy.

After metadata screening, the procedure selects earliest, middle and latest
complete seven day windows. A complete shorter daily window is used in
the broader screening procedure only when source duration or resolution
prevents a seven day choice. The nine retained windows are all seven days,
with both interval endpoints retained. This gives 2017 timestamps at five
minute resolution or 673 timestamps at fifteen minute resolution. EPRI
timestamps use the source Coordinated Universal Time (UTC) field. ACEN's source-local time is retained
without assigning an undocumented timezone.

Table~\ref{tab:windows} identifies all nine windows and their source time conventions.
\begingroup
\begin{longtable}{@{}p{0.24\linewidth}p{0.25\linewidth}p{0.11\linewidth}p{0.15\linewidth}p{0.18\linewidth}@{}}
\caption{Operational windows and distinct recorded physical states.}\label{tab:windows}\\
\toprule Source and window & Dates (inclusive endpoints) & Timestamps & Distinct physical tuples & Resolution / time basis\\
\midrule\bottomrule
ACEN DUID1 EARLIEST & 2026-05-01 to 2026-05-08 & 2017 & 1065 & 5 min; source local\\
ACEN DUID1 LATEST & 2026-08-25 to 2026-09-01 & 2017 & 1824 & 5 min; source local\\
ACEN DUID1 MIDDLE & 2026-07-02 to 2026-07-09 & 2017 & 1705 & 5 min; source local\\
ACEN DUID2 EARLIEST & 2026-05-01 to 2026-05-08 & 2017 & 306 & 5 min; source local\\
ACEN DUID2 LATEST & 2026-08-25 to 2026-09-01 & 2017 & 645 & 5 min; source local\\
ACEN DUID2 MIDDLE & 2026-07-02 to 2026-07-09 & 2017 & 994 & 5 min; source local\\
EPRI EARLIEST & 2024-01-02 to 2024-01-09 & 673 & 421 & 15 min; UTC\\
EPRI LATEST & 2024-12-24 to 2024-12-31 & 673 & 324 & 15 min; UTC\\
EPRI MIDDLE & 2024-06-24 to 2024-07-01 & 673 & 445 & 15 min; UTC\\
\end{longtable}

\endgroup


Exact rational serialization of the reported SOC values determines repeated
states. The physical tuple comprises source, device, SOC and available SOC
limits. There are 7729 globally distinct tuples among 14121 timestamps, with
no exact tuple shared between the retained windows under this definition.
Measured power is a contextual time series field rather than a tuple
coordinate. A separate engineering bin width of 0.1 percentage point is
used only for descriptive SOC coverage. Binning does not replace exact
state identities or establish statistical independence.

\subsection{Eligibility for a complete native query assessment}
Applying the criterion requires more than a visited SOC. The native model
must identify its state and directional flow definitions, measurement side,
sign convention, time base, coefficients, operating constraints and any
controller rules that determine admissible operation. An energy scale is
required when constraints are expressed in stored energy. Auxiliary
consumption, retention and multiport relationships must be specified when
they enter the model. An economic objective is not required for this
objective independent device comparison.

For ACEN, registered directional efficiencies of 0.9 and 0.9 are available;
missing efficiency is not the final eligibility limitation. The native
energy tracker, dynamic bid and ancillary service constraints still lack
the required complete link to BMS SOC and power. For EPRI, the documented
5000 versus 5140 kWh capacity discrepancy remains unresolved. An inverter
efficiency is not substituted for battery directional efficiencies.
Neither measured source supplies a complete native model for the same
query set certification used in the simulations.

The broader source screening also distinguishes usable physical data from
a complete model. StoreNet requires a verified household/controller identity
and consistent grid alternating current and photovoltaic direct current
port definitions; parameters from a related publication are not attached
automatically to an unlinked trace. Yangon data include timestamped SOC
and directional powers, but reported mean round trip performance does not
identify two directional efficiencies, and inaccessible energy management
rules and operating thresholds are not estimated from observed extrema.
Access blocked source archives remain access limitations, not nonexistent
datasets or negative physical validation results.

Sign and alignment diagnostics retain gaps, interpolation qualifications,
state resets and observed excursions. They do not clamp SOC to make a trace
fit a candidate model. A unit efficiency energy balance diagnostic does not
estimate efficiency, and an integrated state column is not automatically
a measured state. Native model incompleteness is distinct from an empty
feasible set or a solver timeout. No complete native query was assessed
for any of the nine measured windows.


Panel (a) retains ACEN dashboard AC power in MW; panel (b) retains EPRI battery DC power in kW. These points describe visited physical conditions and do not identify native capability curves or certified query sets.

FIGINSERT4END




\section{Data and software interfaces}\label{sec:reproducibility}
The accompanying computational package supplies the saved inputs, classifications,
exact vectors and implementation sources supporting this document. Table~\ref{tab:data_interfaces}
identifies the main entry files. The complete file descriptions, source units
and missing value conventions are given in \texttt{DATA\_GUIDE.md}.
Fraction strings preserve exact values, and decimal conversion is used only
for display. Records are linked by case and instance identifiers.
Expanded source descriptions retain the complete earlier source, device and
operational metadata fields. Device records are linked by their stable anchor
names; an earlier display ordinal is not a new device identifier. The data
guide also distinguishes historical route settings from the horizon study
environment.

\begin{longtable}{@{}p{.43\linewidth}p{.51\linewidth}@{}}
\caption{Entry points to the accompanying computational files.}\label{tab:data_interfaces}\\
\toprule File within the package & Scientific function\\\midrule
\bottomrule
\texttt{data/Internal\_Pairs.csv} &468 before/after pairs with fixed parameter and eligibility links\\
\texttt{data/Internal\_Assessments.csv} &936 scientific state slots; the 72 regression slots are separate\\
\texttt{data/Internal\_Saved\_States.jsonl} &Complete saved state results and exact native witness vectors\\
\texttt{data/External\_Instances.csv} &All 34 executed device parameter instances, including source units and terminal rules\\
\texttt{data/External\_Saved\_States.jsonl} &Frozen inputs and saved results for 3772 external states\\
\texttt{data/Source\_Mapping\_Details.json} &Source equations, retained rows, transformations, exclusions and unavailable fields\\
\texttt{data/Repeated\_Timing\_Summaries.csv} &All 108 input/route summaries, including noncompleted repetitions\\
\texttt{data/Directional\_Objective\_Values.csv} &Exact values for 13 states and 8 directions, including all 24 displayed entries\\
\texttt{data/Temporal\_Query\_Classifications.csv} &The 65 recorded state and query classifications\\
\texttt{data/Exact\_Hull\_Constructions.json} &Nine endpoint, midpoint and reconstruction records\\
\texttt{data/Boundary\_Quantified\_Results.json} &Twelve saved quantified domain checks; grid records are separate\\
\texttt{data/Simulation\_Saved\_States.jsonl} &Saved inputs, certificates and exact witnesses for 86 complete simulated states\\
\texttt{data/Measured\_Window\_Summaries.csv} &Nine measured windows and physical state counts; no native query outcomes\\
\texttt{verify\_saved\_records.py} &Portable verification of saved records using exact arithmetic; no optimization is run\\
\end{longtable}

The default entry can be run from the computational package without access
to the author's original working directory. It checks counting units and
saved native feasibility, query preservation and reconstructed violations.
It preserves the accepted equality certificates without rerunning their
logical decision. Numerical reproduction of a scientific solve is distinct
from verification of a saved result. The original model and route functions,
recorded library versions and source parameter inputs identify that separate
computational procedure. The companion retains their original module names
and relative hierarchy under \texttt{code/original\_layout/}. Its dependency
guide separates model inspection, saved record verification and scientific
solver execution. Historical campaign drivers require their documented
configuration and output layout; they are not run by the default verifier. The terminal adapter retains the two fixed building
energy floors before any native verification.

Source driven simulation files preserve the calibration rule, recorded
trajectories and input fingerprint. The original market price response must
be obtained from its provider under the applicable terms; it is not included
in the review export. Physical measurement summaries instead retain
their original source, time base and eligibility limits. The original ACEN
and EPRI raw measurement archives and copyrighted literature PDFs are not
redistributed as author generated data. The documented provider references
are \url{https://acenrenewables.com.au/new-england-bess/} and
\url{https://msites.epri.com/period}. Access to the original material remains
subject to the provider's terms. These files have not been assigned a public
repository DOI.

The editable Word source, bibliography and native figure projects accompany the PDF. Document rendering is separate from scientific solver execution.
The corresponding data guide and machine readable files provide the
reproducibility interface for each scientific section.


Earlier algebraic and implementation checks are retained in
\texttt{historical\_validation/}, with a scope guide and original record
identities. These include affine row identities, positive dynamics, model
scope refusals and restricted operating envelopes. Their designs are
separate from the horizon and literature panels reported here. Historical
within interval switching records use a different reference trajectory
space. Neither they nor earlier implementation tests add assessments to
the current numerical totals.

\subsection{Additional certificate and system interfaces}
The later certificates are stored under \texttt{extensions/certification/}.
Their identity table links each originally unresolved input to its accepted
certificate and the unchanged paired design. System input tables, exact
query certificates and saved request records are supplied under
\texttt{extensions/synthetic/} and \texttt{extensions/rts/}. The extension guide
identifies the original implementation hierarchy and verification entry
points. These files supplement the original data package; they do not
replace its classifications or timing records.


\section{Additional certificates for the fixed horizon design}
\label{sec:sat_certificate_followup}

\subsection{Population and certificate types}
The original panel contains 936 assessments in 468 paired settings. Its
recorded outcomes remain 277 equality certificates, 331 strict inclusion
certificates, 152 unresolved assessments and 176 inapplicable assessments.
The subsequent certification addresses exactly those 152 unresolved input identities.
No horizon, initial state, coefficient, native constraint or query is
changed. Table~\ref{tab:sat_certificate_types} separates the proof routes.

\begin{table}[htbp]
\caption{Additional certificates on the original unresolved inputs.}
\label{tab:sat_certificate_types}
\begin{tabularx}{\linewidth}{@{}p{.27\linewidth}rX@{}}\toprule
Certificate & Inputs & Acceptance condition\\\midrule
Structural preservation & 104 & Complete native row preservation and an exactly feasible point; all are lossless instances of the declared net or gross power model.\\
Rational row upper bounds & 45 & Every required reconstructed row and positive coefficient branch is covered; 3,417 exact nonpositive upper bounds are verified.\\
Rational violating queries & 3 & Exact relaxed feasibility, complete query preservation, exclusivity and dynamics, and a strictly positive reconstructed native residual.\\\bottomrule
\end{tabularx}
\end{table}



\begin{table}[htbp]

\caption{Original and cumulative classifications by horizon.}
\label{tab:si_horizon_followup}

\begin{tabular}{@{}rrrrrrr@{}}
\toprule
 & & \multicolumn{3}{c}{Original protocol} & \multicolumn{2}{c}{With later certificates}\\

$T$ & Applicable & EQ & STRICT & Unresolved & EQ & STRICT\\
\midrule
4 & 64 & 44 & 20 & 0 & 44 & 20\\
5 & 64 & 42 & 22 & 0 & 42 & 22\\
6 & 60 & 38 & 22 & 0 & 38 & 22\\
7 & 60 & 36 & 24 & 0 & 36 & 24\\
8 & 60 & 36 & 24 & 0 & 36 & 24\\
10 & 60 & 36 & 24 & 0 & 36 & 24\\
12 & 52 & 30 & 22 & 0 & 30 & 22\\
14 & 52 & 15 & 22 & 15 & 30 & 22\\
16 & 48 & 0 & 20 & 28 & 28 & 20\\
20 & 48 & 0 & 20 & 28 & 28 & 20\\
24 & 48 & 0 & 22 & 26 & 26 & 22\\
32 & 24 & 0 & 12 & 12 & 12 & 12\\
48 & 24 & 0 & 16 & 8 & 8 & 16\\
96 & 24 & 0 & 16 & 8 & 8 & 16\\
128 & 24 & 0 & 16 & 8 & 8 & 16\\
192 & 24 & 0 & 14 & 10 & 8 & 16\\
256 & 24 & 0 & 15 & 9 & 8 & 16\\
\midrule
Total & 760 & 277 & 331 & 152 & 426 & 334\\
\bottomrule
\end{tabular}

\smallskip
\begin{minipage}{\linewidth}
The horizon $T$ counts operating intervals. Both classifications concern
the same applicable inputs. The 176 inapplicable assessments remain outside
these counts. Later certificates resolve the original unresolved inputs
without changing parameters, constraints or queries. They do not change
the original computation times or supply additional experiments.
\end{minipage}
\end{table}



The structural route requires the complete native model. Losslessness alone
does not imply preservation of arbitrary additional constraints. For the
retained rows, cancellation preserves stored energy, reduces both
nonnegative directional powers and preserves their difference. Separate
ratings, the shared rating and the declared net or gross state dependent
power rows are then preserved. The saved nonempty point is checked against
the same native equations.

The row upper bound route uses the reduced formulation in
Eq.~\eqref{eq:reduced_lp}. Its rational multipliers have the required signs.
Any remaining stationarity residual is bounded using valid variable bounds,
and its correction is included in the final upper bound. Equality is
accepted only after all required branches have a nonpositive bound and
nonemptiness is established. A nonpositive floating optimum is not used as
the certificate. The three strict cases have explicit rational witnesses.
The 45 upper bound cases extend through 32 intervals; equality at longer
horizons in this subsequent certification comes from the structural route. The three
new strict cases occur at 192 and 256 intervals.

\subsection{Versioned state and pair accounting}
With the additional certificates, the same applicable population contains
426 equality and 334 strict inclusion assessments. The original
classifications describe the original protocol and budget; they are not
overwritten by the later certificate record. Table~\ref{tab:sat_pair_accounting}
shows the corresponding pair accounting.

\begin{table}[htbp]
\caption{Pair classifications before and after additional certification.}
\label{tab:sat_pair_accounting}
\begin{tabular}{@{}lrr@{}}\toprule
Pair outcome & Original & With additional certificates\\\midrule
Equality to equality & 106 & 172\\
Equality to strict inclusion & 64 & 82\\
Strict inclusion to strict inclusion & 123 & 126\\
At least one unresolved assessment & 87 & 0\\
Inapplicable update & 88 & 88\\\midrule
Total predefined pairs & 468 & 468\\\bottomrule
\end{tabular}
\end{table}

The 18 newly classified transitions are members of the original paired
design. Their feasible charging updates preserve the original horizon,
efficiencies and native constraints. They are not an additional sample of
18 experiments. The saved energy feasible power violation witnesses for
the original 64 transitions remain a separately defined witness population.
Neither the additional certificates nor repeated timing executions increase
the number of predefined inputs.

\subsection{Timing evidence and reproducibility}
The original 36 input repetition study in Section~\ref{sec:computation}
and the later matched route comparison have different purposes. The latter
uses six original inputs, two routes and three repetitions, giving 36 timing
executions. Both routes use the same processor affinity, one linear algebra
thread, alternating order and a 30~s primary certification budget. The
archived route remains unresolved in all 18 executions, while the new route
returns verified certificates in all 18. These are censored completion
comparisons, not a finite speedup ratio. The original route also attempts
an additional state of charge constrained witness query within its budget.
The comparison concerns the full certification workflows, rather than
isolated solver kernels. It establishes no general complexity bound or
real time capability.

The accompanying state reconciliation identifies every original input,
original outcome, additional certificate and pair membership. Certificate
checking requires no optimization: it checks structural row conditions,
nonempty points, rational multipliers and bounds, or saved violating queries.
Historical solver outputs and timings retain their own protocol versions.



\section{Selected requests in the synthetic system case}
\label{sec:sat_selected_requests}

\subsection{Declared model and scheduling outputs}
The illustrative case combines the public synthetic three bus network and
load profiles with study specified storage and service constraints.
The source is a pinned PyPSA example snapshot\cite{pypsa_examples_snapshot_23e2436}.
The predefined settings, profile identities and source digests are retained
in the registered system design. The network is an illustrative model,
rather than a measured grid. Four one hour intervals are used. The storage
energy range is 0 to 6~MWh, its shared power limit is 1~MW, and the main
lossy settings use charging efficiency 0.9, discharging efficiency 0.95 and
retention 0.99. Optional state dependent charging limits, lossless and
other context controls are specified in the registered settings. Such
capability constraints are study assumptions, rather than capability curves
measured by the public network source.

The three generators each have a 100~MW upper bound and zero lower bound,
with cost coefficients 10, 20 and 100 in the source objective units.
The lines joining buses 1--2, 2--3 and 1--3 have respective ratings
100, 100 and 10~MW. Each normalized direct current (DC) flow is the difference
of its endpoint angle coordinates; bus 3 is the angle reference. Load is
placed at bus 3, and storage at bus 1. The declared angle bounds are
$[-1000,1000]$ in model coordinates. The profiles and their exact load
fractions are supplied in \texttt{REGISTERED\_SYSTEM\_CASES.json}.
No physical currency or field calibration is inferred from these normalized
network and objective coefficients.
The optional taper is the end of interval native row
\begin{equation}
-p_t+\beta E_t\leq 1\ \mathrm{MW},\qquad
\beta=\frac{1}{6}\ \mathrm{h}^{-1}.
\label{eq:sat_synthetic_taper}
\end{equation}
This is distinct from the steeper charging limits in the controlled
state transition study. The representative Figure~9 request uses the box
device, without this optional taper row.

Thirteen registered paired settings give 26 state assessments. Their
before and after members are related by the declared feasible first action.
They are not 26 statistically independent networks. The query retains net
power and every generator output, branch flow and bus angle throughout the
horizon. The reference adds charge and discharge exclusivity to the same
native model. Reconstruction fixes the exclusive powers from net power and
recomputes stored energy, leaving all retained variables unchanged.

Economic dispatch minimizes generation cost. For the service request,
the sum of net power over the last two intervals is minimized under a
common cost cap equal to 1.01 times the exclusive economic optimum at
that state. Cost is then minimized at the optimal primary value.
This cap is shared by the relaxed and exclusive service problems at a
given state, but can change between members of a state pair because their
economic optima differ. It is a request selection condition and does not
restrict the complete global query set used for equality assessment.
The saved service value uses a positive benefit convention in the data
table; the corresponding minimized objective value is its negative.
Agreement of this primary value does not establish agreement of the
secondary costs. The original selection and tie rules are unchanged.


\subsection{Public load transformation and registered settings}
The variable profile uses the London column of the public meshed example
load file. Its first four observations are divided by the maximum over
the complete ten observation source column and multiplied by 100 MW.
London is a source column label; this transformation supplies an illustrative
load profile rather than metered system demand. The resulting loads are
3.664552010121, 100, 25.653256677878 and 13.385529985646 MW.
Their exact fractions and the separate flat 100 MW control are retained in
the registered input file.

The source network notebook supplies the triangle topology, normalized
reactances, line ratings and generator cost coefficients. Placement of the
storage, its energy and power ratings, loss allocation, retention, free
terminal energy and optional taper define the combined study model.
Auxiliary consumption is zero. The seven pilot settings were registered
before solving; the six context settings were admitted by the registered
pilot outcome rule. No new settings were selected for this document.
Table~\ref{tab:registered_system_settings} lists every retained pair.
The updated initial energy follows the first exclusive economic action.
The same fixed load profile and complete horizon are then reassessed;
this operation is not a rolling horizon shift.
\begin{table}[htbp]
\caption{Registered synthetic system settings.}\label{tab:registered_system_settings}
\begin{tabular}{lllllll}
\toprule
Pair & Profile & Device & $\eta_c/\eta_d$ & Initial (MWh) & Updated (MWh) & Relation\\\midrule
1 & Variable & Box & 0.90/0.95 & 6/5 & 643/4750 & EQ to EQ\\
2 & Variable & Box & 0.90/0.95 & 3 & 3643/1900 & STRICT to EQ\\
3 & Variable & Box & 0.90/0.95 & 24/5 & 8786/2375 & STRICT to STRICT\\
4 & Variable & Taper & 0.90/0.95 & 6/5 & 643/4750 & STRICT to STRICT\\
5 & Variable & Taper & 0.90/0.95 & 3 & 3643/1900 & STRICT to STRICT\\
6 & Variable & Taper & 0.90/0.95 & 24/5 & 8786/2375 & STRICT to STRICT\\
7 & Variable & Box & 1/1 & 3 & 197/100 & EQ to EQ\\
8 & Flat & Box & 0.90/0.95 & 6/5 & 261/125 & EQ to EQ\\
9 & Flat & Box & 0.90/0.95 & 3 & 387/100 & STRICT to STRICT\\
10 & Flat & Box & 0.90/0.95 & 24/5 & 1413/250 & STRICT to STRICT\\
11 & Flat & Taper & 0.90/0.95 & 6/5 & 1044/575 & STRICT to STRICT\\
12 & Flat & Taper & 0.90/0.95 & 3 & 387/115 & STRICT to STRICT\\
13 & Flat & Taper & 0.90/0.95 & 24/5 & 2826/575 & STRICT to STRICT\\
\bottomrule\end{tabular}\end{table}

The pair numbers in this table follow the registered case order. Complete
identifiers, exact load fractions and query vectors are provided in
\texttt{extensions/synthetic/REGISTERED\_SYSTEM\_CASES.json} and
\texttt{extensions/synthetic/SYSTEM\_STATE\_PAIRS.csv}. Numerical selection used
the recorded 60 s solver limits. Exact feasibility and request preservation
were checked separately in the saved certificate implementation.

\subsection{Counts and the representative service request}
Table~\ref{tab:sat_selected_counts} separates complete query classifications
from the selected request outcomes. A strict query inclusion certificate
does not imply that every selected request fails.

\begin{table}[htbp]
\caption{Selected request outcomes in the 26 predefined assessments.}
\label{tab:sat_selected_counts}
\begin{tabular}{@{}lrr@{}}\toprule
Complete query relation & Economic request & Service request\\\midrule
Equality: 7 states & 7 realizable & 7 realizable\\
Strict inclusion: 16 states & 16 unrealizable & 16 unrealizable\\
Strict inclusion: 1 state & 1 realizable & 1 unrealizable\\
Strict inclusion: 2 states & 2 realizable & 2 realizable\\\midrule
Unrealizable selected requests & 16/26 & 17/26\\
States with equal primary values & 10 & 14\\
Unrealizable within that subset & 0/10 & 5/14\\\bottomrule
\end{tabular}
\end{table}

The representative request is the selected service output for the
\texttt{PUBLIC\_FIRST\_FOUR\_BOX\_E4of5} before state, with initial energy
4.8~MWh. Table~\ref{tab:sat_service_energy} gives its stored energy and
native upper bound margin. The final relaxed energy is exactly 6~MWh.
The exclusive final energy and its excess are
\begin{equation}
E_4^*=\frac{1326208859289}{220281250000}\ \mathrm{MWh},
\qquad
E_4^*-6\ \mathrm{MWh}=\frac{4521359289}{220281250}\ \mathrm{kWh}.
\label{eq:sat_service_final}
\end{equation}
The right hand quantity is 20.52539328245141~kWh. It is the service request
violation; the economic request at the same state has a feasible exclusive
reconstruction. The primary service benefit is 2 in both model optimizations,
while the secondary costs differ. The exact complete query is preserved
by the reconstruction that violates the bound.

\begin{table}[htbp]
\caption{Representative service request energy and native bound margin.}
\label{tab:sat_service_energy}
\begin{tabular}{@{}rrrrr@{}}\toprule
Boundary & Relaxed (MWh) & Exclusive (MWh) & Headroom (kWh) & Margin (kWh)\\\midrule
0 & 4.800000 & 4.800000 & 1200.000000 & 1200.000000\\
1 & 3.699368 & 3.699368 & 2300.631579 & 2300.631579\\
2 & 4.294460 & 4.315402 & 1705.540251 & 1684.598109\\
3 & 5.151515 & 5.172248 & 848.484848 & 827.752128\\
4 & 6.000000 & 6.020525 & 0.000000 & $-20.525393$\\\bottomrule
\end{tabular}
\end{table}

Headroom is the native upper bound minus relaxed energy; margin is that
bound minus reconstructed energy. Positive reconstruction shifts at earlier
boundaries remain feasible because headroom is still available. Exact
fractions, complete request vectors, every native residual and the full
26 state mapping accompany the source data; displayed decimals are rounded.
The Figure~9 display uses an explicitly labeled inverse hyperbolic sine
scale for the signed margin and retains the full energy context. Connecting
lines join saved boundary values and do not certify an intermediate
violation time.



\section{Supporting network certificates and request scope}
\label{sec:sat_network_scope}

\subsection{Mapped benchmark and unchanged query}
The Reliability Test System developed by the Grid Modernization Laboratory
Consortium (RTS GMLC) supplies the fixed 73 bus, 120 branch network and
24 hourly inputs for 1 January 2020. The source is the pinned public RTS GMLC dataset\cite{rts_gmlc_source_3ece0d3}.
The mapped continuous direct current (DC) dispatch model includes generator
capacity and ramp constraints, nodal balances, branch flow equations and
ratings, reference fixed bus angles and the declared storage equations.
It uses zero generator lower bounds and does not reproduce unit commitment,
reserve, reactive power or a complete original plant controller. The
source head and tail volumes are represented by the study's equivalent
affine energy model. Shared power and optional charging taper constraints
are study assumptions. Verification concerns this complete declared model.

The global query retains all 24 storage net powers and all generator, flow
and angle coordinates. Economic and service budgets used to select requests
do not restrict that global query domain. A policy constructing a network
extension may choose generator outputs, but a reconstruction of a given
complete query must preserve its already specified outputs.


\subsection{Network and storage parameter mapping}
Source bus 101 fixes the reference angle. The DC branch coefficient is
the 100 MVA base divided by source reactance and tap ratio; a zero source
tap denotes a ratio of one. Loads are allocated from the hourly regional
series in proportion to the source bus demands. The continuous dispatch
retains 153 active generators. The three synchronous condensers and the
unmapped concentrating solar unit are excluded; source storage at bus 313
is represented separately.

Wind, photovoltaic, rooftop photovoltaic and hydro availability is the
smaller of source nameplate capacity and the hourly series. Thermal
marginal cost uses source fuel price times the first incremental heat rate,
divided by 1000, plus variable operation and maintenance cost. The average
heat rate at zero output is the documented fallback. Hourly ramp limits
are 60 times the source per minute rate. Generator lower bounds are zero,
and no preceding horizon ramp is imposed. These choices define a
continuous dispatch comparison, rather than a unit commitment reproduction.

Source storage has 50 MW charging and discharging ratings, 0.15 GWh upper
and 0.075 GWh initial head volume. The equivalent head energy range is
0 to 150 MWh, initially 75 MWh; tail energy is 150 MWh minus head energy.
The study allocates the source 85 percent round trip efficiency to charging,
with discharging efficiency one. The lossless control sets both to one.
Retention is one, auxiliary consumption is zero and terminal energy is free.
The shared rating and optional end of interval gross charging row are
declared study constraints:
\begin{equation}
c_t+d_t\leq 50~\mathrm{MW},\qquad
c_t+\frac{1}{3}~\mathrm{h}^{-1}E_t\leq50~\mathrm{MW}.
\label{eq:rts_study_rows}
\end{equation}
The second row is included only in taper settings. The energy coordinate
is an affine equivalent of source reservoir volume; it is not a calibrated
battery capability model. Reactive power, voltage magnitudes, AC losses,
commitment, reserves, contingencies and the complete plant controller are
outside this declared dispatch model.

The RTS service objective maximizes net discharge over the six hours with
the highest gross load, with earlier hours breaking ties. These are source
periods 10, 17, 18, 19, 20 and 21. This differs from the last two interval
charging request in the synthetic case. Its cost cap uses 1.01 times the
numerical exclusive economic optimum. Secondary cost selection uses a
primary objective band of $10^{-6}$, rather than exact lexicographic
optimization. The original numerical solve limit is 30 s, within a
240 s state limit. Native numerical checks use absolute tolerance $10^{-5}$
and relative tolerance $10^{-8}$; cost agreement uses the larger of 0.01
and $10^{-8}$ times the absolute exclusive cost. Exact set certificates
remain separate from these numerical request checks.
\begin{table}[htbp]
\caption{Complete query certificate inputs for the mapped RTS case.}\label{tab:rts_certificate_inputs}
\begin{tabular}{lllll}
\toprule
Setting & Member & Initial (MWh) & Relation & Native rows\\\midrule
L0 / T0 & new & 25 & EQ & 25038\\
L0 / T0 & old & 75 & EQ & 25038\\
L0 / T1 & new & 375/4 & EQ & 25062\\
L0 / T1 & old & 75 & EQ & 25062\\
L1 / T0 & new & 75 & STRICT & 25038\\
L1 / T0 & old & 75 & STRICT & 25038\\
L1 / T1 & new & 7050/77 & STRICT & 25062\\
L1 / T1 & old & 75 & STRICT & 25062\\
\bottomrule\end{tabular}\end{table}

L0 and L1 denote lossless and lossy conversion, while T0 and T1 denote
absence and presence of the study taper. The two identical 75 MWh members
of L1/T0 are two assessment slots at one parameter state. The eight records
therefore cover seven distinct complete states. Native row counts include
both directions of equalities; they are not numbers of independent tests.

\subsection{Network extension and complete query certificates}
The saved exact network baseline sets the colocated generator to
177.5~MW. Every storage power trajectory in the declared power box has a
network extension under
\begin{equation}
g_{313,t}=177.5-p_t,\qquad -50\leq p_t\leq50.
\label{eq:sat_network_policy}
\end{equation}
All other generator outputs, flows and angles remain at their exact baseline
values. Colocation cancels the net injection change. The generator range
127.5 to 227.5~MW respects its 355~MW capacity; its greatest possible
adjacent change is 100~MW, below the 248.4~MW ramp limit. These checks cover
the continuous power box, rather than a finite set of sampled trajectories.
The saved row certificate checks all 24,798 network rows, including
4,656 equality rows; a separate implementation checks the network relations.
The policy establishes network feasibility, rather than storage feasibility
for every power trajectory in the box.

When the device and network couple only through net power and use the same
network constraints, a device query without an exclusive realization
remains a missing complete system query whenever it has a feasible network
extension. Conversely, preservation of the device query preserves those
network rows because the reconstruction leaves all net powers and retained
network variables unchanged. This is a consequence of the shared model
structure, rather than a new general network theorem. Constraints directly
coupling network decisions to internal energy or directional powers require
the complete native constraint test again.

Four original and four updated assessments have complete query certificates:
four establish equality in the lossless controls and four establish strict
inclusion in the lossy models. They represent seven distinct parameter
states because one update has zero first action. Equality certificates
check every native preservation condition and exact nonemptiness. Strict
certificates provide exactly feasible complete relaxed queries and their
unique, query preserving, infeasible exclusive reconstructions. A
restricted sign search used to construct a violating witness is not used
as a proof of equality.

\subsection{Selected numerical requests and nearby exact points}
The original eight selected network requests are two request types for each
of four original model states. They passed numerical native reconstruction
checks. A separate panel contains six saved optimization outputs per model
state: economic, primary service and secondary service outputs from each
member of the model pair, giving 24 outputs. These outputs are not 24
additional network model states.

All 24 outputs have saved nearby rational feasible points, whose exclusive
reconstructions also pass the native model. Comparison with the original
saved decimal queries gives identical net power in two cases, identical
retained network coordinates in none, and an identical complete query in
none. Exact nearby feasibility therefore does not certify exact preservation
of the original requested outputs or exact optimality of their floating
solutions. The nearby conditions retain the original tolerances: $10^{-5}$
for powers and energies, $10^{-7}$ for angles, a cost allowance of the larger
of 0.01 and $10^{-8}$ times the original absolute cost, and a primary service
face tolerance of $10^{-6}$.

The recovery comparison uses the same four model states and 24 outputs.
Four exclusive economic outputs were development inputs; 20 were held out
at the output level. They are not independent held out model states. The
generic exact nearby implementation timed out on all 24 outputs under its
specified budget; the recorded physical elimination and rational candidate
route accepted all 24 after exact checking. Its selected active faces were
zero dimensional for this panel. This finite comparison supports neither
completeness of the recovery heuristic nor a general solver speedup claim.
All timeouts and failed diagnostics remain in the research archive.

The network extension also explains a negative engineering result: binding
line limits in selected outputs do not show that the network excludes
storage queries, because the feasible colocated policy leaves baseline
flows unchanged. This fixed case establishes complete certificate closure
and distinct request guarantees; it does not establish congestion caused
failures, field performance or a broadly representative economic effect.
